% =====================================================================
%  R. Nielsen, OM AANDSDANNELSE med særligt Hensyn paa vort Opdragelsesvæsen.
%  Sex Forelæsninger.  Kjøbenhavn: Gyldendal, 1877.  109 pp.
%  / On Spiritual Formation, with Special Regard to Our System of
%    Upbringing.  Six Lectures.
%
%  TRANSLATION (English).  Translated 2026-10-06 from transcription.tex (Danish) in
%  this directory, the text of record, as committed at 1b28f4f (collated 2026-10-06,
%  OPEN 0; see its header).  Mirrors it paragraph for paragraph.  \opage{N} at the English word where printed page N begins, placed by
%  ebook-method/collate/pagemarks.py from the Danish markers (anchors in
%  ebook-method/om-aandsdannelse/anchors-en.txt).  Method: ../../../TRANSLATION-PLAYBOOK.md.
%  The 55 odd spots in the Danish noticed while translating (.parts/tr/oddspots.tsv) were
%  read at the image the same day and the transcription corrected through its pipeline
%  (see the transcription's header, TRANSLATION READING); this English was then brought
%  into line: p. 12 „forlorent Dogme“ (no dash), letterspacing on pp. 7, 31, 60, the bold
%  „først“ on p. 66, and the comments at the sites revised.
%
%  CONVENTIONS.  „…“ -> ``…'' (quotations inside quotations `…').  Letterspacing
%  (\emph) carried over one for one; Latin \textit as printed.  Titles Nielsen gives in
%  quotation marks keep them; titles he letterspaces keep \emph.  "d. e." / "ɔ:" -> "i.e.";
%  "f. Ex." -> "e.g.".  Quotations from Kierkegaard, Ørsted and Scripture are translated
%  from Nielsen's Danish (he quotes Kierkegaard closely, with his own cuts and
%  punctuation).  Shelley (p. 41) is given in Nielsen's Danish prose and translated from
%  it, with Shelley's lines in a comment.  Misprints (PRINTED AS IS in the Danish, and
%  the unmarked ones found while translating) are translated by their sense, with a
%  comment at the site.
%  TERMS (as in the repo's other Nielsen translations, esp. om-hindringer):
%    Aand / aandelig = spirit / spiritual · Dannelse = formation (dannet = cultivated)
%    Aandsdannelse = spiritual formation · Kulturdannelse = cultural formation
%    Kulturmenneske = man of culture · Opdragelse = upbringing
%    Opdragelsesvæsen = educational system · Underviisning = instruction, teaching
%    det Mangfoldige / det Enfoldige = the manifold / the simple · Erkjendelse = cognition
%    (Verdens-, Gudserkjendelse = cognition of the world, of God) · Viden = knowledge
%    Tro = faith · Autoritetstro = faith in authority · Tilegnelse = appropriation
%    Selvforstaaelse = self-understanding · Inderlighed = inwardness · Forargelse = offense
%    Existensmeddelelse = existence-communication · Gjennembrud = breakthrough
%    Forbillede = pattern / exemplar · det Humane / Humaniteten = the humane / humanity
%    Nutiden = the present age · Kirkeskole / Kulturskole = church-school / culture-school
% =====================================================================
\documentclass[12pt,a4paper,openany]{book}
\usepackage[utf8]{inputenc}
\usepackage[T1]{fontenc}
\usepackage{libertinus}
\usepackage{libertinust1math}
\usepackage{textalpha}
\usepackage[english]{babel}
\usepackage[protrusion=true,expansion=true]{microtype}
\usepackage{enumitem}
\usepackage{textcomp}
\usepackage{setspace}
\usepackage[margin=1.3in]{geometry}
\usepackage{fancyhdr}
\usepackage{hyperref}
\hypersetup{pdftitle={On Spiritual Formation (1877)}, pdfauthor={Rasmus Nielsen}, hidelinks}
\pagestyle{fancy}
\fancyhf{}
\fancyhead[LE,RO]{\thepage}
\fancyhead[C]{\textit{R. Nielsen: On Spiritual Formation (1877)}}
\setstretch{1.3}
\usepackage{marginnote}
\newcommand{\tocline}[3]{\noindent\makebox[2.2em][r]{#1}\hspace{0.6em}#2\dotfill\ #3\par}
\newcommand{\opage}[1]{\marginnote{\footnotesize\textit{[#1]}}}
% a lecture opens a new page: heading, a short rule (as in transcription.tex)
\newcommand{\lecture}[1]{\clearpage\begin{center}{\large #1}\\[0.4em]\rule{3em}{0.4pt}\end{center}\medskip\noindent}

\begin{document}

% title page (p. [I]); the publisher's device is not reproduced
\begin{titlepage}
\begin{center}
\vspace*{3em}
{\Huge On Spiritual Formation}\\[1em]
{\large with Special Regard to Our System of Upbringing.}\\[2.5em]
\emph{Six Lectures}\\[1em]
by\\[1em]
{\large \emph{R. Nielsen.}}\\[5em]
\emph{Copenhagen.}\\[0.5em]
Gyldendal's Bookshop, Publishers (F. Hegel \& Son).\\[0.5em]
{\small Printed by J. Jørgensen \& Co.}\\[0.5em]
{\small 1877.}\\[4em]
{\small Translated from the Danish}
\end{center}
\end{titlepage}
% p. [III], Contents
\begin{center}{\large\textbf{Contents.}}\\[0.3em]\rule{3em}{0.4pt}\end{center}
\hfill{\small Page.}\par
{\small
\tocline{}{\emph{First Lecture}: The Difference between Cultural and Spiritual Formation.}{1}
\tocline{}{\emph{Second Lecture}: Faith in Authority and Personal Self-Understanding}{17}
\tocline{}{\emph{Third Lecture}: Spirit and Letter: The Critical Turning Point}{38}% as printed: the lecture begins on p. 34
\tocline{}{\emph{Fourth Lecture}: S. Kierkegaard's Appearance}{52}
\tocline{}{\emph{Fifth Lecture}: S. Kierkegaard's Self-Contradiction}{71}
\tocline{}{\emph{Sixth Lecture}: Religion and Humanity}{91}
}

\lecture{First Lecture.}
\opage{1}Great and considerable efforts are being made in our day for the benefit
of enlightenment and formation. Alongside learned formation and the
manifold kinds of professional formation, weight is laid on general
formation: political, aesthetic, social formation. Everyone in society
takes pride in counting as a cultivated person. It is with formation as
with dress. Formation is a becoming adornment; who does not wish to be
well dressed and to show himself to advantage?

But despite the commendable zeal for the many-sidedness of formation, it
might well appear, if we look more closely, that something has been
forgotten, that the much-praised many-sidedness is in a high degree
one-sided. What the enlightened of the present age understand by
formation is as good as exclusively a formation in relation to
\emph{the manifold}; it is what we will call, in one word, \emph{cultural
formation}. That there must also be a formation in relation to the
simple, and that this is the proper \emph{spiritual formation}, is, so
far as we can see, all too lightly overlooked. It is assumed without
further ado that cultural formation is all in all, so that when one only
has culture in ample measure, spiritual formation follows of itself; that
is a dubious error.

\opage{2}Cultural formation turns its gaze outward; whatever glides past the
searching eye in the world of multiplicity is transformed into
phenomenon. Thought, the supersensible, is phenomenon just as much as the
sensible; conscience is phenomenon just as much as electricity, genius a
phenomenon as much as lightning, Christ a phenomenon as Buddha and
Mohammed are. If we now call the human being whose whole content of
consciousness empties itself out into the forms of cultural formation a
man of culture, it is easy to understand that such a man, however busy he
may be with things and with himself, can never become so turned inward
that he in earnest catches sight of the essence either of things or of
himself. He by no means denies that there is essence in things, but, he
says, this essence in and for itself cannot be known. He does not deny
that in a way he does have a soul, but what sort of being the soul is,
and whether it really is a being---how should he be able to know that? If
he will contemplate his soul in order thereby to discover his own self,
then this self instantly becomes a phenomenon to him. And what, then,
lies at the ground of this phenomenon? Something mystical, over which one
should not brood. Indifferent to what he is in himself, he then
cultivates himself so as to seem something, to be seen and to cut a
figure as a considerable phenomenon in the series of phenomena.

Turned outward toward the manifold, the cultural consciousness is
many-sided; it is like an Argus studded with eyes, a hundred clear eyes;
but these hundred clear eyes all lie under the spell that whatever they
fasten upon instantly hides its essence and turns into mere phenomenon.
But one cannot immerse oneself in a phenomenon that keeps its essence
hidden, and \opage{3}that in which one cannot immerse oneself has only a fleeting,
passing interest. One grows tired of familiar phenomena, as the child of
its toys, and craves something new. No single phenomenon has validity in
and for itself; only the alternation of phenomena, their play of colors,
and their connection in large groups have interest. Once one is finished
with them at home, one must be off and look about abroad. The man of
culture is a cosmopolitan, an Odysseus, widely traveled; he speaks many
languages, has been in many lands, has lived with Cyclopes and nymphs
alike. How enlivening it is to hear him speak of the mountains he has
climbed, the museums he has visited, the excavations of priceless
treasures of antiquity at which he himself was present. He can criticize
authors, for he knows the literatures; he can describe buildings and
monuments, for he knows the styles; paintings and statues!---yes, he has
a connoisseur's eye, he has taste.

But what does the man of culture know in relation to the simple? What a
question! If simplicity is a phenomenon like everything else in the
world, should the cultivated man not know that besides stupidity there
are higher kinds of simplicity: in the innocence of the child, in the
virginal soul of woman, in the artist's naivety of genius, in the faith
of the common people, in religion? Yes indeed, he says, holy simplicity
is noble and beautiful; but criticism is the flower of formation, the
consummation of culture; criticism, ``the tenth Muse,'' stands to
simplicity as the falcon to the dove; one poisonous sting of sharp
criticism can kill faith, and with faith religion falls.

Spiritual formation turns its gaze inward, yet not so as to exclude the
manifold. In optics there are lenses of different shapes, some that
\opage{4}scatter the rays, others that gather them. Cultural formation prefers
ray-scattering lenses with an imaginary focal point; spiritual formation
requires gatherers of rays and places them so that the focal point falls
in the self. ``Know thyself!''---theoretically, that means: fix your
attention on your inner, supersensible being, for only in and through
this being are you able to know phenomena. It is, after all, you yourself
who see and hear, perceive and feel; it is you yourself who rejoice and
grieve, think and ponder; it is you yourself who live. If you would know
what life is, what the soul is, what consciousness is, then seek the
answer above all in yourself. The stars of heaven, the animals of the
earth and the atoms of the brain cannot give you the least contribution
at first hand, for what you know of these things you know only by virtue
of your own consciousness; what you judge of these things you judge by
virtue of your own understanding; what value you attach to these things
you attach to them by virtue of your own feeling. Do not imagine that by
starting from a critical consideration of the animal species in their
multiplicity, or of the human race in general, you can discover what life
is in its essence; critical consideration must have a standard, but
whatever you discuss and whatever you measure, in the last instance you
carry the standard within yourself. ``Know thyself!''---practically, that
means: hold fast to yourself, not in selfish conceit but in free,
conscious interaction with the surrounding world. There is an order of
things under which you must bow and by which you must let yourself be
determined. It prescribes conditions to you, but you can choose; it
offers you enjoyments, it lays sufferings upon you, it threatens you with
death; but \opage{5}how you understand all this depends on your composure of mind
and your self-understanding; how you conduct yourself in all this depends
on the direction of your will, and so, essentially, on how you yourself
are.---From dispersion in the external, from the multiplicity of
phenomena, spiritual formation would lead the human being back to the
simple One, the essential inner being in which the self lives, and then
through this inner being to the source of all life, the author of nature,
to the riddle of necessity, to the wellspring of freedom, to spirit as
spirit, to God.

Without spiritual formation no coherent world-view, cast in one piece, is
possible; without a coherent world-view no true independence of character
is possible. An inwardly free and independent character can no more be
acquired by sailing with every current and keeping a lookout for every
curiosity than by shutting one's eye to the actual and, like a
navel-gazer, staring at the mystical point in the soul. What we call, in
the essential sense, a world-view, a view of life, rests on a reciprocal
determination of what comes from without by what is originally in
ourselves and comes from within---a reciprocal determination so
thoroughgoing that the soul, amid all the impressions from the world of
things and phenomena, preserves the impression of itself. The self must
show itself as a co-determining, fundamentally determining factor, not
merely as the collector and bearer of the whole. The essence of this
many-sided reciprocal determination, rooted in self-determination, is
spirit.

Spiritual formation is not attained by accepting on external authority
the fixed dogmas of a positive religious doctrine. There are many who
have plenty of \opage{6}positive dogmas about spirit, and yet they have no
spiritual formation; whereas a ``free thinker,'' who denies all dogmas,
can have spiritual formation. As regards the negative and the positive in
spiritual formation, it depends essentially on how the denial and the
acceptance are grounded. If someone, e.g., will deny the existence of God
and the immortality of the soul because he imagines that sufficient,
irrefutable proofs of the unreasonableness of accepting them have been
given in natural science and criticism, then he is blind and stupid in
relation to spirit. If, on the other hand, he is in a position to
understand himself and to carry his life through consistently in denial;
if he finds that there is something in his being that rebels against God
as well as against eternal life; if he denies because he is clearly
conscious that, by the constitution of his mind, he will deny and must
deny: then his denial bears the stamp of originality; he has spirit in
negative form, he has spiritual formation. A denier of spirit who has
spirit can be an instrument of spirit, an able taskmaster for a
generation that is spiritually enfeebled, spiritually unclear,
spiritually lazy, spiritually befooled.

Positive spiritual formation can go in different directions, according
as one stands toward the two ideal powers: \emph{necessity and freedom}.
If necessity is posited as the original and freedom as the derived,
formation will go in a rational direction, and conviction will be
grounded on knowledge. If freedom is posited as the original and
necessity as the derived, formation will go in a religious direction, and
conviction will be built on faith. In the first case the emphasis falls
on the cognition of the eternal reason that expresses itself in nature;
the human being grasps and understands himself as \opage{7}a rational being in a
world of rational, necessary laws. In the second case the emphasis falls
on the cognition of God and the word of life that reveals his will. The
human being grasps and understands himself as an ``image of God,'' an
organ of will for the divine spirit.

But in whichever of these two main directions spiritual formation goes,
the cognition of the world must still allow itself to be reconciled with
the cognition of God, and the view of life must be coherent in both
respects. From this standpoint it can best be seen how matters stand with
the many-sided formation of the nineteenth century, namely by casting a
glance at our educational system.

In the doctrine of reason one has got no further than a critical
dissolution of the dogmas of positive religion, and in the doctrine of
religion no further than a simple repetition and modified inculcation of
the old dogmas. If upbringing were to be built on a view of life in which
there really was unity and coherence, one would have to have, at least
for the time being, two separate institutions of upbringing, two kinds of
school: a \emph{culture-school} on the foundation of reason and a
\emph{church-school} on the foundation of positive religion, so that
there really would be two views of life to choose between.

In the church-school everything would then have to be strictly kept away
that could mediately or immediately weaken faith in the historical
infallibility of the religious tradition and awaken doubt about the
absolute truth of the dogmas. The whole of upbringing would have to be
arranged accordingly. The church-school could then not confine itself to
instruction in religion alone; it would have to give, in its own way, a
real cognition of nature and the world, normative for culture and
science. The pupils in such a school \opage{8}might well learn languages, both
ancient and modern; but for the origin of the many languages they must
know no other explanation than the building of the tower of Babel. That
history must be taught, both Bible history and general world history,
goes without saying; but the decisive thing is that the course of history
and the great events be kept constantly in the light of Paradise and the
Fall and the devil's continued struggle against governing Providence. Nor
is there then anything to prevent mathematics, physics and natural
history from being taught as well; but it must count as an unshakable
presupposition that the world was created in six days, that heaven is a
fixed vault, and that the earth stands still while the sun rises and
sets. Hypotheses about animal species that are supposed to have died out
before the appearance of man, about a natural state of savagery without a
preceding Paradise, and the like, must be treated as ``the rudiments of
the world,'' fantasies of an unbaptized reason; for it is surely clear
that no animal could die before death had come into the world, and death
came, after all, only after man had sinned. In this way unity comes into
instruction. Man knows how he was created, how he has fallen, and how he
can be raised up again: on this foundation an upbringing can be carried
through, and spiritual formation will be accordingly.

In the culture-school the ground must be cleared and the premises
cleansed of all manner of musty traditions. Criticism, with ``free
thought,'' should see to it that the place is properly aired, so that one
can breathe freely in the atmosphere of reason. In place of the dogmas of
theology the great principles of natural science are set up: the
eternity of matter, the immutability of the laws of nature. Astronomy's
teaching about the globes in \opage{9}space, geology's account of the periods of
the earth, zoology's conception of the descent of the human race, the
general presupposition of nature as the absolutely necessary,
self-sufficient ground of everything, spirit included: these cognitions
of reason furnish the elements of the foundation on which the view of
life is to rest and upbringing is to be laid out. A moral teaching is
necessary; but the moral commandments must be conceived as general
commandments of reason, whose fulfillment is demanded by society, but
whose application within the limits of the coercive laws receives
freedom and freshness, color and life, from the peculiar natural ground
of the self-conscious individuality. Religion can be taught just as well
as cultural history and languages, but the whole development falls under
the viewpoint of the necessity of reason; it goes without saying that the
positive religions, with their miracles and supernatural phenomena, are
without exception conceived as mythologies. Add to this a demonstration
of the ideal in art and poetry, an acknowledgment of the nobler feelings
of the human heart, of the right of conscience and of the free
personality: and one surely has a coherent world-view on which a system
of upbringing can be carried through consistently. Spiritual formation
will be accordingly.

Let it be well noted that the two systems of upbringing are so worlds
apart in spirit and direction that any principle borrowed from the one
and taken up into the other must necessarily confuse the mind and weaken
the character. Now, has one hitherto kept the two schools apart from each
other? No. Is it likely that one will be able to keep them apart in the
future? No. What, then, has one done, and what will one \opage{10}go on doing? Mix
them together in total unclarity. From the common school right up to the
university the mixture is continued, in our critical present age, in a
most uncritical manner. One thinks thus: a cultivated person must know
the rational order of the world, and a cultivated person ought to know
the chief doctrines of religion. Rational cognition guides us in our
undertakings; religion strengthens our moral feeling and comforts us in
our sufferings. Society can do no more for the young generation than open
to it access to many-sided formation; the relation between reason and
positive religion is a riddle that each must solve for himself as well as
he can.

But is it now also defensible to fill young souls with a hodgepodge of
conflicting notions and to ground their upbringing on a sum of dark
riddles that one neither can nor cares to solve oneself? Let religion
have its mysteries and nature its own; it is nonetheless an inescapable
demand that those who are to be the teachers of youth must be able to
give an account of how matters stand with our cognition of the
supernatural in religion and with our rational cognition of the world.
If they cannot, then humane upbringing in the present age, as regards the
formation of character, must come to stand far below the Catholic
upbringing of the Middle Ages and the Greek and Roman upbringing of
antiquity.

Imagine among us a father's address to his son when the son, his
instruction completed, leaves school. ``My son! You now know the two
guiding stars that are to shine upon your way, religion and reason. Never
lose sight of them! They seldom shine equally brightly \opage{11}at the same time,
but when the one goes out, the other is lit, so you will never be without
guidance. The one sits high in the supernatural heaven and gives you
hints of an almighty will which, when it will, can shake the earth and
make the stars fall; the other shines in the infinite space of nature and
makes clear to you a world order so unchangeable that not even the least
speck can be moved in it without a natural cause. The one rises in the
east, the other in the west. By following the first you are taught that
the earth is under a curse, and that the way to the goal is the way of
sufferings to blessedness in the world beyond; by following the second
you know that nature is good as it is, and this world an arena for a rich
development of your gifts and powers. Religion has taught you that before
God you are a guilty sinner, eternally lost if you cannot be pardoned
through another's merit and intercession. Reason has let you see that you
must yourself atone for your transgressions, that no substitute and no
arbitrary grace can help you out of the contradictions of life in which
you have entangled yourself. So I know that you have been many-sidedly
instructed. Steer your course now by the guiding stars! With your double
compass you can safely plow the sea of life and expect to reach a quiet
harbor.''

Such an address will sound like a parody of our whole educational system;
and yet it is so free of all exaggeration that the contradictions
indicated become far more glaring when one goes into particulars. But it
is one thing to mix the irreconcilable, another to combine what
essentially belongs together. If one grants that a religious cognition of
God is just as necessary to man as a rational cognition of the world, and
\opage{12}that spiritual formation is conditioned by the union of the two, then the
world of culture has a great task to solve with regard to its system of
upbringing.

In the teaching of the special scientific disciplines there is nothing in
this respect to correct. The choice of material and method depends on
the purpose and on particular conditions. Everything depends on the
application that the results of science find in the doctrine of spirit.
We choose as an example \emph{creation}. Is this world created, or is its
nature, with absolute necessity, what it is from eternity? To judge by the
haste with which modern criticism has decked out positive religion as a
sacrificial victim on the altar of naturalism, one would certainly
suppose that creation was a sham dogma, a blind notion without rational
meaning. But is it really so? Natural science starts from the principle
that there is reason in nature. That is more than a hypothesis, for if
the opposite is assumed, science is impossible. If someone should perhaps
assert that reason and nature are one and the same, so that the
expression ``there is reason in nature'' says nothing other than that
there is reason in reason, nature in nature, then such an assertion does
not exactly betray an excessive measure of critical acumen. If there is
to be rational meaning in the expression, it must necessarily be
understood that the unity of reason and nature includes an essential
difference. If we will now stop for a moment and reflect on what is then
properly to be understood by reason, might not logical necessity lead us
to spirit?

``\emph{The laws of nature},'' says H.~C. Ørsted, ``together with what is
active, constitute the \emph{universal} in \opage{13}multiplicity and the eternal in
change. The laws of nature can therefore also be called \emph{thoughts of
nature}. But just as in all thinking there is an activity, and the
thought, thought without it, is only a form of thinking, so in every
existing thing there is an activity by which it asserts its existence and
expresses it against other things. Without activity it has no actuality.
In thinking we call this active element will, in external objects force.
But just as willing and thinking are one, only regarded from two sides,
so too are force and natural law; they are expressions of an infinite
unity of force and thought, of divine willing and divine
\emph{thinking}\dots{} What we know in nature is thus nothing foreign,
but the same as what lives in ourselves.''---If, then, from ``what lives
in ourselves'' we are to be able to see what divine thinking and willing
mean, it must surely become clear to us that, even if we cannot possibly
picture the Godhead to ourselves as a thinking and willing individual in
human fashion, we must nevertheless necessarily conceive it as a
thinking, willing self-being, i.e., as a rational being distinct from
nature. An eternal thinking without an eternal thinker, an eternal
willing without an eternal willer, thoughts that no one has thought,
determinations of will that no one has willed---such assumptions belong
in the world of fancies. Nature can no more think the thoughts that
express themselves in its laws than it can will the ``actions'' by which
these laws are fulfilled. Every thing acts according to its nature; but
what is a thing's nature other than its inner disposition, determined by
divine thinking and willing? Since no natural thing can give itself a
disposition, \opage{14}the rational disposition of nature must originally be
determined by the all-determining self-being, i.e., by the all-knowing,
all-powerful spirit. Through the cognition of things and their laws we
are led to the cognition of a rational world order; through the cognition
of a given world order to the cognition of the ordering reason; but the
ordering reason is a fancy if it does not have spirit as its ground and
bearer.

With the same right as we say ``there is reason in nature,'' we can also
say ``there is reason in spirit.'' Spirit and reason are one, but in such
a way that this unity includes a difference. When we speak of reason, we
think first and foremost of the infinite content of thoughts of wisdom,
determinations of will, concepts and laws that is realized in nature by
virtue of the dispositions of things. When we speak of spirit, we think
of the infinite self-being which, thinking, willing, producing, unfolds
the content of reason from its inner depth. From this standpoint the
relation between necessity and freedom, nature and creation, can be
cleared up. Necessity is seen in nature in so far as it reveals itself in
the rational connection of things: the consequent must necessarily have
its ground, the effect its cause. To assume an arbitrary will which,
without natural causes, could take it into its head to produce or
annihilate natural things, or to make natural things unnatural, is
certainly a sign of dreaming and superstition.

``It has indeed also been attempted,'' says Ørsted, ``to defend the
dreams of superstition against natural science by pointing out that an
immeasurably greater part of nature is more unknown to us than well
known. This defense can, \opage{15}even among the uninitiated, show nothing but a
possibility, which must nonetheless be all the less significant the more
of the old superstition's fancies have been so illuminated that their
nothingness has been settled; but for those who have not been content to
appropriate natural science with memory and understanding, but have lived
their way into its whole, it is clear that our insights, despite all
their imperfections, nevertheless annihilate every thought that anything
contrary to reason should take place in nature.'' We must unconditionally
endorse this statement.

Within the given whole, things and laws must necessarily correspond to
each other: other forces, other laws; other laws, other forces. But is
the absolute necessity of the given whole thereby proved? By no means! To
assume that the given world of experience must be the only possible world
because every breach of its laws is contrary to reason is just as
unreasonable as to assume that the right-angled triangle is the only
possible triangle because the denial of the Pythagorean theorem is
contrary to reason. The given nature is certainly a rational whole, but
the peculiar character of this whole is determined by its peculiar
disposition, originally determined by an almighty self-determination of
spirit as spirit, and so by a \emph{creation}.

If, now, from the standpoint of the doctrine of reason the prospect has
been opened that the rational cognition of the world can be united with a
religious cognition of God, then it must surely also be possible to show,
from the standpoint of the doctrine of religion, that the religious
cognition of God is compatible with a rational cognition of the world.
But on what conditions? Here there is something that must be changed,
\opage{16}remolded, reworked from the ground up. What is it? Surely not revealed
religion, its supernatural origin, its miracles and holy facts? Far from
it! Rationalism, the new as well as the old, has prostituted itself more
than enough in this respect. It is not the world-view of revelation, as
it stands before us with all the marks of antiquity, whose features of
character are to be effaced. No, it is something quite different that
must be demanded for the sake of upbringing and spiritual formation. It
is our inner, essential relation to what has been revealed, our
understanding, our personal appropriation of it---in short, our whole
religious consciousness, our faith---that must be reworked from the
ground up. What such a reworking means, where it leads, and how far it
can be carried through, will be the subject of closer consideration in
what follows.

\lecture{Second Lecture.}
\opage{17}Religion is the foundation of the life of the spirit; religion must be
grasped in faith. It is then natural to make a person's faith the
measure of his spirit: the stronger the faith, the more vigorous the
spirit; the more vigorous the spirit, the stronger the faith. Yet this
is not always so. True, there can be no true and fruitful development of
the spirit without faith; but there can in actuality be an obstinate,
strong and unshakable faith that is nonetheless very poor in spirit, and
indeed incapable of grounding any real spiritual formation.

For faith has two main sides, just as positive religion has two main
sides: an external and an internal, a sensible and a supersensible, a
historical in its tradition and a suprahistorical in its essence. In the
beginning the two elements are fused together: the internal comes to
light in something external, and the external converts itself
immediately into something internal; but gradually they separate, and
dangerous crises set in in the life of the spirit.

It is a characteristic mark of all positive religions without exception
that at their first appearance they announce themselves as divine
revelations. The less spirit there is in such a revelation, the \opage{18}more
emphasis is laid on the acceptance of the external. The orthodox
Mohammedan assumes that the Koran has existed from eternity in the
highest heaven on the so-called Preserved Tablet, and that from there it
was revealed piecemeal to Mohammed by the angel Gabriel over a period of
23 years. The contradictions found in the holy book are resolved by the
interpreters in such a way that the verse revealed later abrogates the
earlier one. Here is a firm, unshakable faith in something external.
Judaism, too, lays decisive weight on what has externally happened and
what is externally prescribed; but, as telling proof that there is
originally more spirit in the faith in Jehovah than in Mohammedanism,
the prophets stood up against the ceremonies of the priests and the
external commandments, and in the name of the spirit made the internal
prevail. In Christianity, the absolute religion of the spirit, the
external must serve the internal. Christ says: ``my words are spirit and
are life.''---``I am the truth.''---``You shall know the truth, and the
truth shall make you free.''---If one gives up the positive, historical
foothold, inwardness becomes mystical, and the religious consciousness
hovers in the air; but if, conversely, one places the assurance of faith
in the sensible experience of a supernatural external, then the
Mohammedan faith stands, in life and vigor and in glowing colors, on the
highest step.

In the history of the Crusades wonderful things are told of a wonderful
figure in the East, ``the \emph{Old Man of the Mountain}.''
Hassan-ben-Sabbah, head of the ``Brotherhood of the Ismailites,'' had
known how to instill in his subjects so unconditional a faith in the
words of the Koran, so punctual an obedience to its commandments, and
thereby so blind a submission to his \opage{19}own holy authority, that he
astonished the world. His faithful despised worldly honor, suffering and
torments, indeed life itself, in order to obey his slightest sign. That
was upbringing. A Count Henry of Champagne traveled through the land of
the Ismailites, visited the Old Man in his castle, and received an
honorable welcome. On each of the high towers two men stood on guard,
rigid as pillars. Hassan gave a sign to a couple of these sentries, and
instantly they hurled themselves from the tower and lay shattered at the
feet of the astonished crusader. The Old Man remarked in cold blood:
``Should you wish it, you would see them all hurl themselves to the
ground like that at my sign.''

With this unconditional faith, this blind obedience, there was, according
to the accounts, a peculiar connection. So far were the men on guard on
the towers from facing such a moment with horror, or regarding the Old
Man as a cruel tyrant obeyed only out of fear, that on the contrary they
honored him in the inmost depths of their souls and regarded a relief
from duty such as the one described above as the happiest of all. The
closer circumstances are presented thus. The sect of the Ismailites was
originally divided into two classes, masters and apprentices; to these
Hassan added a third class, which was to be ignorant of the sect's real
secrets. The members of this class were called the Victims; they were
dressed in white, with red caps and sashes, and had to keep constantly
about the Grand Master in order to defend and avenge him. In the midst of
the lands over which the Old Man ruled lay a castle, a paradise of
Mohammed on earth. Lovely gardens, leafy trees, flowers bedewed by the
spray of fountains bade the one who entered welcome; shady walks shielded
him \opage{20}from the burning of the sun. The dwelling within shone with gold and
marble; arabesques, silks and costly carpets in the high, cool halls,
young girls resting on soft divans in voluptuous luxury!---here was
everything a believing soul could wish and desire. Now when a youth was
to be initiated for admission to the class of the Victims (the
Fedawi), it was done in the following manner. In the Old Man's presence
he was stupefied by a drink prepared from hashish (the leaves of Indian
hemp), and in this unconscious state he was carried to the enchanting
place. When he then opened his eyes, he could not but believe that he
was in Paradise, and enjoyed in full measure all the joys of
blessedness. He was then stupefied again and brought back. The Old Man
stood once more at his side and assured him, now awake, that he had not
bodily been away from the spot; if, notwithstanding, it was as he said,
then his soul must by a miracle have been caught up into Paradise. So the
young man had faith in his hands; in melancholy he longed for the hour
when he should enter there for all eternity, and he was willing to submit
in unconditional obedience to the conditions that the holy authority
prescribed for him.

Even within the religion of the spirit there is something Mohammedan in
the faith of the church, especially the Catholic church. What we here
understand by the Mohammedan element does not lie in the first instance
in the teaching, in the dogmas, in the precepts, in the ceremonies, not
in what is the object of faith, but originally in faith itself, in the
way in which faith is conceived and determined. If one should perhaps
object that the way in which the essence of faith is conceived and
determined is itself a dogma, then one is merely moving in a circle. One
cannot possibly prescribe faith in the \opage{21}dogma of faith without
presupposing a personal inner being, a conscience, that can believe
itself obliged to accept the dogma. However much persuasion or coercion
one employs, the acceptance cannot become faith unless there is
something in the soul itself that says Yes and Amen to the dogma of
faith. Faith is therefore, in the first instance, an original
expression of freedom, an act of self-determination, a recognition of
truth on personal grounds. In what, then, does the Mohammedan element
consist? Simply in this: that the personal impression of the object of
faith, which should form the starting point for an appropriation in
spirit and in truth, is all at once transformed into a motive for
unconditional submission to an external authority. ``You must surely
feel within yourself and recognize that you are a fallen, helpless
creature in need of restoration and salvation'': when this appeal to the
conscience of man has had its effect, then follows the summons: ``seize
the salvation that is offered, as the shipwrecked man seizes the plank
of rescue; let yourself be received into the bosom of the one Church
that saves!'' Now it is all over with self-understanding and free
appropriation. Faith in the Church, blind submission to its authority,
unconditional obedience to its commandments is the condition of
salvation. To be sure, the Church does not neglect to inform its faithful
of what it is they are to believe. It teaches them what God is, what
holiness is, what sin is; all this can be clearly understood down to
the last detail. But from all this understanding there is one thing
essentially excluded, and that is self-understanding. It is the Church's
God, the Church's holiness, that is believed in. The believer does not
know at all what God is in himself; he knows only what the Church teaches
that God is. He does not even know what sin is and what it is to sin;
\opage{22}the Church must literally teach him that. What the Church stamps as sin
is sin, even if natural feeling says the opposite; when the Church has
blotted out a sin, then it is blotted out, however much the original
sense of right and wrong may have to object. The secret is that the
believer has no conscience of his own, that he has given the Church his
conscience. Faith in authority is externality in the internal; the
formula of this externality is the dogma.

The alarm has been sounded against the Catholic Church's many historical
encroachments in external matters; and it is true, they have been
outrageous, and may perhaps yet, when ``the modern consciousness'' has
properly helped along and done its best, become outrageous enough; but
these are the ordinary encroachments of the wielders of power, which are
met and halted with external weapons. The radical encroachment, the one
that devours the spirit, is the transformation of the internal itself
into something external. This Mohammedan element---that faith in
authority takes the inner man captive, lures him out of himself, makes
him a stranger to himself, so that, stupefied as by a magic potion,
fantastically and yet with understanding, he can find blessedness
outside himself and thereby attain a kind of sensible assurance of its
reality---is a beguiling illusion, an insurmountable barrier to
spiritual formation.

Much outcry has been raised over the many arbitrary additions in the
Catholic tradition, its invented miracles, its legends and superstition;
not without reason. The miracles are its supernatural building
materials, its ornaments; the legends twine like ivy up its walls: it is
luxuriously furnished with the inventions of superstition. And yet the
objects of superstition are not yet the essence of superstition, but
only its \opage{23}outgrowths. Superstition sits in the sanctuary, in faith itself
as faith in authority.

For all that, there is spirit in the Catholic way of being. On the
foundation of authority and blind obedience the Jesuits have shown what
can be attained by mystically compressed inwardness and a consistent
system of upbringing. A character like Francis Xavier one will seek in
vain in any other church communion. Such courage of faith, such
perseverance in renunciation and self-denial, such concern to win souls,
must, in the eyes of those who look one-sidedly at the external, place
this Jesuit missionary on a level with the apostle Paul. But according to
his inner being, what is he then? An incarnation of that faith in
authority which robs the personal self of its power from the ground up.
It is not his own human being that has come to development under the
influence of the Church; on the contrary, it is the being of the Church
that has forced its way into him from without. The Church has taken his
heart, his will, out of his breast and set its own heart, its spirit,
its will in their place. He believes the Church's faith in the Church's
name; his renunciation of the world is in him the Church's renunciation,
his zeal for conversion the Church's zeal; he himself is the blind
instrument, ``will-less as a corpse,'' obedient ``as the staff in an old
man's hand.'' What he does is certainly something far higher and better
than what those men did who hurled themselves down from the tower,
because what his authority has prescribed to him is higher and better;
but the motive is the same: unconditional obedience, blind submission to
an external authority.

Through its faith in authority and its system of upbringing the Catholic
Church has set an eternal enmity between itself and human reason. It may
indeed seem \opage{24}as if the systematic development and grounding of the
Church's body of doctrine contained an essential acknowledgment of the
right of reason; but the acknowledgment is only a semblance. Certainly
there have been, and still are, great men of science, famous scholars,
not only historians and philologists but also mathematicians,
astronomers, physicists, and even philosophers, in the bosom of the
Church. But if we look more closely, the sciences that move in a worldly
domain are only tolerated to a certain degree. The Church can, as
circumstances dictate, close an eye to scientific theories and
hypotheses so long as they do not openly set themselves against its
authority; but this always happens on the tacit presupposition that its
own infallible knowledge is infinitely exalted above ``the rudiments of
this world.'' As regards thinking in relation to dogma, the dependence of
reason is downright slavish. Reason is obliged to concede that Christ has
two wills because the Church says he has two; if the Church said he had
two and a half wills, % the print has "kvis" for "hvis" (PRINTED AS IS)
then reason would not only be compelled to accept it as truth, but would
even have to do its best to get it proved. That the bread which, before
it is consecrated by the priest, is natural bread can after the
consecration keep its natural taste, its natural appearance and its
natural properties, and yet no longer be natural bread, but be
transformed into the body with which Christ was crucified---this is so
mighty a blow in the face of reason that everything must go black before
its eyes and convince it that the Catholic faith in authority is just as
incompatible with a rational cognition of the world as the revelations
of the Koran.

\opage{25}With the Reformation the dogma of faith began to undergo an essential
transformation. Luther broke with the Catholic authority, asserted the
right of conscience, and conceived faith as the personal appropriation of
the merit of Christ, and so as the fundamental factor in the inwardness
of the religious life. For a moment it looked as though the outer was to
be reduced all at once to stuff and material for the inner as it worked
its way forward, and the holy words of Scripture, glowing through with
spirit, kindled the light of self-understanding in the soul. But it was
not long before the object of faith drew attention away from the
movements in the believing self. The strong religious life became little
by little a doctrine about life. The Protestant Church got its own
theology and its own dogmatics in the likeness of the Catholic:
externality celebrated a new triumph when the Church became a state
church and orthodoxy went so far as to have its dogmas registered, like
deeds, by the police authority. Under royal custody, faith in authority
had now taken up residence in Scripture, essentially just as external as
formerly in the papacy.

The assertion that something is true because it stands in the Bible is
just as Mohammedan as that something is true because it stands in the
Koran. Orthodox faith in the Bible can no more catch sight of the circle
in which it moves than orthodox faith in the pope. He who believes in the
infallibility of the Church imagines that, freed from all human
fallibility, he has the infallible over against him, and does not notice
that it is he himself who adds the infallibility out of his own. He
believes in the Church because it is infallible, and behold, then it is
only his own faith in its infallibility that is infallible. Likewise with
him who builds his faith externally on Scripture. He \opage{26}would not know at
all what he was to believe if the infallible Scripture did not teach it
to him; but neither does he know that it is his own---infallible!---faith
that teaches him that Scripture is infallible. So on both sides one
appeals to the testimony of the spirit; but with faith in authority one
gets only the dogma about the spirit, not the spirit itself, only the
dogma about the testimony, not the testimony itself. In the personal
testimony of the spirit the divine must be confirmed by the human just
as fully as the human by the divine.

It is a fact that Lutheran orthodoxy, held fast in its positive dogmas,
stands in the same hostile relation to the culture and science of the
present age as the Catholic. It is not only in the physical and the
historical, but even in the psychological, in the life of
consciousness, that the obstacles to a mutual understanding show
themselves. Think of the doctrine of the absolute corruption of human
nature through original sin, of God's universal grace, which is
nevertheless not universal but benefits only the elect; of the doctrine
of the two natures in Christ, of his seat at the right hand of the
Father, of the omnipresence of his body and yet its limited presence in
the Lord's Supper, etc.; and one will soon discover, without overtaxing
one's understanding, that even if many, many have believed literally in
all this, there is nevertheless no one who has believed it in spirit and
in truth, for the simple reason that no one has been able to believe it.
To believe in something is to have it in the form of dogma as the
object of faith; but to believe something is, by inward appropriation,
to transform it into the content of faith. As regards objects of faith,
there can hardly be anything shown so unreasonable and absurd that it
cannot, from one side or another, become \opage{27}the object of a faith in
authority: the absurd is, after all, precisely that which can be
accepted only on authority. It is otherwise with the content of faith.
What is to be inwardly taken up into the religious consciousness as
spiritual nourishment for the soul must, like every other nourishment,
let itself be dissolved and digested; it must assimilate itself to the
essence of consciousness: the understanding of it must become one with a
personal self-understanding. The positive dogmas, be they for the rest
knotty and angular or round and polished, are in their fixed state like
stones, clear crystals, perhaps precious stones!---but not the bread of
life. One can hang them on the outside of consciousness and adorn one's
soul with them; but one cannot digest them.

The relation between the school of reason and the church-school appears
in a new light when faith is conceived as the personal appropriation of
the universally religious content of the positive dogmas. It is by no
means the task to abolish the faith of the church, or in any external
sense whatsoever to attack authority as authority. On the contrary! In
external relations authority can be not only venerable but even
necessary. Not even rational instruction can do without authority. The
disciple must accept definite propositions on the teacher's authority,
and, trusting that what still lies above his understanding will become
intelligible to him, work his way forward to insight. The coat of arms of
authority ought to stand over the entrance to the school of inquiry. A
science that has not yet won esteem in public opinion, and thereby a
certain authority, will be exposed to the most superficial objections
without anyone troubling to look into the matter. If it goes thus with
the rational, which must after all be able to assert itself by virtue of
universally valid \opage{28}proofs, how must it not go with the religious, which
can be recognized only on personal grounds? If the word of the Bible has
no authority, the Church and its dogmas no venerability, no esteem, then
religion is abandoned, and anyone can mock it according to whim and
fancy. The truth itself, the eternal truth, which is before Abraham was,
had, after all, to put on flesh and blood and step forth before men in
human form as authority.

But even the absolutely true authority is only something external; the
truth, the eternal being, is something internal; the truth demands
recognition, and for the sake of recognition it reduces the external to
a means and gives its authority over to dissolution. ``You shall know the
truth, and the truth shall make you free.'' That authority is dissolved
does not, then, mean that it is abolished; no, it means that it is in
truth confirmed. Only the false authority is finished when it is
dissolved. A false authority, theology for example, must be weakened the
more, the more its scientific character is tested, for every attempt to
know its grounds dissolves into the absurd. Perhaps it will be objected
that when theology loses its authority, the Church and the dogmas must
suffer for it. That would be true if the authority of the Church were
built on theology; the rational thaw that dissolves the last remnants of
theological authority would then certainly draw the dissolution of the
Church's authority after it. But the authority of the Church, when we
think about it, is built not on theology but on religion; it is an
authority not for science but for faith, an authority of faith. Whether
this authority of faith is true or false depends on whether its dogmas,
through religious, personal appropriation, dissolve \opage{29}into truth or into
untruth. Here is a task for spiritual formation, a fundamental problem
for a religious-rational upbringing. Some examples.

\emph{Nature and creation}. ``In the beginning God created the heaven and
the earth'': these words of the Bible, with the story of creation that
follows them, have the authority of the Church. ``Nature is uncreated;
all formations in heaven and on earth have come about naturally'': this
assertion has, according to the naturalistic way of seeing, scientific
authority. Can an unconditional faith in either of these two authorities
now take the place of the recognition of truth? Impossible! If one lets
the authority of the Church exclude that of science, one can easily get
free of all trouble with what is rational. One lays one's head to rest
on the lazy pillow of authority, reason dozes off, and reason, gently
dozed off, is not easily disturbed by the noiseless tread of rational
knowledge. % the print has "ikke let at af" (an extra "at")
One dreams of the miracle of creation and Adam in Paradise; but such
dreamings do not lead to a personal recognition of truth in faith. If
one lets the authority of science exclude that of the Church, one
undeniably has a spur to rational inquiry into everything given in its
connection; but whence the connection and the connected are remains a
riddle, for personal self-knowledge is lacking. If religion has no
authority, one is all too quickly finished with freedom in the question
of the origin of things. Nature is necessity and necessity nature;
natural necessity swallows up freedom, first the divine and then the
human. If science has no authority, one is all too quickly finished with
nature. The Creator is almighty and can do everything he will \opage{30}with
created nature; man's freedom is not essentially dependent on the laws
of nature but on the free will of the Creator: the whole dissolves into
divine arbitrariness. Only the united pressure of both authorities can
set all the wheels of inquiry in motion and drive the life of
consciousness forward from the cognition of God to the cognition of the
world, from the religious to the rational, from faith to knowledge and
from knowledge to faith.

\emph{Sin and redemption}. Religion, with the authority of the Church,
will vouch that sin is the ruin of nations, that sinful man cannot by his
own strength rise from his fall, but needs a divine helper and redeemer.
Reason, with the authority of science, will vouch that universal
sinfulness is a simple paraphrase of the necessary imperfection of human
nature, that the conditions of liberation must be sought in the
perfectibility of this nature, that the human race, by a rational use of
its natural powers, can and must in the course of time little by little
become its own savior. If one now throws the authority of reason
overboard, what then? Then the dogmatic notion of sin and corruption
grows to such a height that it darkens the whole of existence and ends
by making religious self-knowledge impossible. No light of faith is able
to illuminate the abyss of corruption. If we conceive faith not as a
faith in dogma, a blind faith in authority, but as a personal
appropriation, grounded in self-understanding, of the revealed word,
then there is surely no one who can or could believe that human nature
has been so corrupted at the root by Adam's sin that the individuals of
the race are without exception forfeit to eternal damnation. If one
rejects the authority of the Church \opage{31}and of dogma, then not only does
original sin become a perverse fancy, but even actual sins find so many
excuses in the drives of natural necessity that sin and guilt become
empty phrases, and the consciousness of sin a sign of effeminate
voluptuousness. Among the spokesmen of culture there may even occur such
superficial, indeed such brutal, scoffing remarks about sin and guilt
that moral accountability is wiped out. Why? Because one does not find
it worth the trouble to go into oneself and examine the consciousness of
sin more closely.---Both authorities, the Church and science, are
indispensable if man, through the inner co-operation of faith and
reason, is to discover the secret of sin and redemption and have his
personal being cleared up.

\emph{The eternal life}. We have the authority of Scripture and of the Church
for the claim that the existence of an eternal life is historically
established by the resurrection of Christ. We have the scientific
authority of natural science for the claim that the story of the
resurrection is in conflict with the laws of nature. Now if the authority
of the Church is to hold exclusively, then the assurance of eternal life
is built on the story of the resurrection, whose authority is built on
the Bible, whose authority is built on the authority of the Church, which
in turn is built on the authority of history. That is a Mohammedan
assurance by virtue of something external. If the authority of
knowledge is to hold without further ado, then the story falls, together
with the dogma, to criticism, and the result is that it is not life but
the idea of life that is the eternal. In both cases one seeks certainty
about the self outside the self, on the Church's side in a history, on
the side of science in a contemplation of nature: the personal
understanding of personal life is excluded. The scientific earnestness
with which \opage{32}the organic conditions of the soul's existence are
investigated must correspond to the earnestness of faith with which
religious inwardness embraces the consummation of personal life, and
vice versa. The inquiry of reason is not to be restricted or adjusted in
a finite manner to the religious, any more than the inquiry of faith is
to be limited and accommodated to the rational. The dissolution must be
complete from both starting points. Only when both authorities are
dissolved, and the dissolution shows that the religious is precisely the
possibility of the rational, and the rational the possibility of the
religious, and that the unity of both possibilities is the original
self-unity of personal life, destined for eternity---only then can faith
be united with knowledge, religion with science.

As the only way out for reconciling the cultural consciousness with the
religious need, people have begun to call back to life the rationalism
that fell asleep more than a generation ago, by breathing the spirit of
the age into it. But the attempt will hardly succeed. The rationalist is
right when he demands that the dogmas of orthodoxy must be dissolved, but
he is wrong to dissolve them into knowledge and not into faith; to
dissolve them into knowledge is to crack them like hollow nuts: the
religious kernel is gone. The rationalist wants to have his articles of
faith grounded in science and does not notice that he smuggles in
something personal that science cannot ground. He presupposes a holy,
righteous and loving God, and does not notice that he has thereby gone
far beyond the limits of knowledge. Concepts such as God's
righteousness, God's love, God's grace, mercy and the like are not
concepts of knowledge but concepts of personality, concepts of faith. The
righteous God of love is a God of personality, \opage{33}a God of the heart, a God
of conscience. The rationalist assumes that the doctrine of the soul's
immortality is rational, at bottom scientific, and does not notice that
he confuses what science cannot disprove with what it can prove.

Lack of insight into the relation between the personal and the
impersonal, the religious and the rational, is lack of spiritual
formation, lack of upbringing.

\lecture{Third Lecture.}
\opage{34}We know from daily life how easily dispute can arise over the relation
between word and meaning. One should not, as the saying goes, cling to
words but look at the meaning; and yet the meaning has its hold in the
word. One should, as it is put in a higher style, consider that the
spirit gives life but the letter kills; but in judicial rulings and in
legislative assemblies one must also consider that the spirit clings to
the literal expression and would without it be homeless and unreliable.
The dispute about word and meaning, spirit and letter, gains decisive
influence in the Protestant Church on the interpretation of Holy
Scripture. The dogmaticians of the 16th and 17th centuries held the
spirit fast in the letter and maintained that God himself had dictated
word for word to the holy authors what they should write. In this sense
the word of Scripture was God's word: prophets, evangelists and apostles
had been only ``God's scribes,'' ``God's pens.'' There cannot, then, be
even the least error or discrepancy in the written document of
revelation, whether as regards the historical facts or as regards the
expression of the religious: Holy Scripture is in each and every thing
\opage{35}literally and absolutely infallible. How theology has since fared with
the deification of the Bible and the absolute faith in its
infallibility, we shall not dwell upon: we pass over the fatalities of
the external and fix our attention on the psychic, the internal.

Without closer acquaintance with the criticism of the biblical writings
and the Church's dogma of inspiration, anyone who combines with the
general formation of the present age a religious need and reverence for
the word of Scripture, and reads through the Old and New Testaments,
will convince himself that the spirit of revelation does not let itself
be caught in the letter. It stands written in the Second Book of Samuel
(24:1): ``The wrath of Jehovah again struck Israel; for he tempted
David''; it stands written in the First Book of Chronicles (21:1):
``Satan was Israel's adversary, and he tempted David to number Israel.''
The event is the same. If one now holds strictly to the letter, then one
gets out of it that Jehovah and Satan are one and the same being. If one
separates Jehovah from Satan, one has by that very act separated the
spirit from the letter, the meaning from the word, and applied an
interpretation. It stands written in the prophet Hosea (11:1): ``Israel
was still young when I came to love him; out of Egypt I called my son.''
It stands written in Matthew (2:15) that the flight of Joseph with mother
and child into Egypt took place ``that it might be fulfilled which the
Lord spoke by the prophet, who says: Out of Egypt I called my son.'' % the print opens this quotation („at det maatte ...“) and does not close it
If this is to be understood literally, then the evangelist is in
contradiction with the prophet, and the spirit in contradiction with
itself. It is neither the same son nor the same event. The son of whom
the prophet speaks is the Jewish people, which in its youth was in
Egyptian bondage. The prophet \opage{36}looks back on what is past, but the
evangelist says that he looks forward to something to come. According to
the letter, then, the spirit, when it literally speaks of the past, was
literally speaking of what is to come. For the sake of the meaning it is
necessary to apply an interpretation.

The interpretation cannot be infallible unless the reader can secure for
himself an infallible interpreter. That the spirit in the word of
Scripture is infallible benefits us only on condition that it can make
our understanding of it infallible. If the reader of Scripture will now
hold that the spirit stands in the very same relation to him as to the
prophets and apostles, and tells him word for word what he is to believe
and do, then his reading of Scripture is, in the religious sense,
superfluous. If he must admit to himself that the spirit worked in the
holy authors in a quite different way from the way it works in him, then
a psychological barrier to understanding is thereby set up. The task is
to determine the right relation between the original and the derived.

People complain of what is doubtful in the biblical narratives, of
contradictions in the historical matter, etc.---all of it something
subordinate. Even if all difficulties were removed, and the Bible from
first to last in literal agreement with itself, it would still not help
us in the very least, so long as the interaction between the divine and
the human spirit remains inexplicable. The wonder of Pentecost is a
telling example. If the outpouring of the spirit is understood
immediately according to the letter, then it can well, in the character
of a miracle, become an object of faith in authority; but how will it
fare with personal appropriation? ``\opage{37}Are not all these who speak
Galileans? And how then do we each hear them in our own tongue, in which
we were born? Parthians and Medes and Elamites, and we who dwell in
Mesopotamia and Judaea and Cappadocia, Pontus and Asia, in Phrygia and
Pamphylia, Egypt and the parts of Libya about Cyrene, and we Romans
sojourning here, Jews and proselytes, Cretans and Arabians---we hear them
speak in our tongues?'' (Acts 2:7 ff.) For faith in authority there is
no difficulty here at all; God is, after all, almighty and can do
everything he will. But for faith as personal appropriation there is a
great psychological difficulty here. If God's spirit, when it will,
intervenes in the life of the soul in so supernatural a way that
suddenly, without any preparation, it imparts to consciousness a
facility in foreign languages, then the miracle of inspiration stands on
the same level as the miracle of Balaam's ass.
The same spirit that can, by supernatural influence, make a human soul
speak a foreign language previously unknown to it can also make an
animal soul speak human language. But what then is the nature of the
soul, and where are we with our rational consciousness? A real miracle
cannot be comprehended, that is true; but if it is not to stand before
us as something wildly foreign, something that does not concern us at
all, then at least its possibility must be recognized, and the
recognition of this possibility must be grounded in self-knowledge.

If we begin from the rational, then one of two assumptions is possible:
either the order of nature must have its original ground in nature
itself, or in a divine willing. If one assumes the first, then there is
no revelation; the doctrine of revelation is mythology. If one assumes
\opage{38}the second, then it must be more closely settled what is to be
understood by a divine willing. One can then set the divine willing in
so necessary a unity with the laws of the sequence of natural phenomena
that it wills nothing and can will nothing except what agrees with our
rational insight into the connection of the whole and can be
incorporated into the ever-widening circles of our experience: from this
rationalist standpoint, too, there can then be no question of any
special revelation of God.

If the possibility of a revelation is to be recognized, then the divine
willing must be expressly placed in a divine will, an almighty,
self-determining freedom, and the moving ground of this freedom in an
infinite love. This conception must renounce scientific proof, but it is
not on that account contrary to reason. No human reason can measure the
depths in the disposition of nature, for no human reason can measure the
depths of divine love. Natural causes rest on natural dispositions,
natural dispositions on divine self-determining; actual nature is a
manifestation of the will of spirit. Torn loose from spirit, nature has
no independence. The natural causes of natural things are all grounded in
the supernatural: the material actuality of nature consists in being the
possibility of the supernatural working of the divine will.

Seen from this standpoint, the universe with its laws can be said to rest
on a wonder. The creating spirit has not spent its freedom on the order
of natural necessity, nor exhausted the inner depth of its being in the
outer world of the universe. If the creating, sustaining will is a \opage{39}will
of love that will reveal itself, then the possibility of a personal
revelation of love must belong from the beginning to the disposition of
nature. The wonder of revelation, with its miracles, is from this side to
be conceived as a sign of God's personal freedom: in the outworking of
the omnipotence of love, the supernatural hidden in nature steps forth
from its hiddenness and shows itself in the natural; something new
appears in nature, but the order of nature is not broken.

What, now, stands in the way of reconciling faith in revelation with a
rational cognition of the world? Among other things, two! In the first
place, the modern fixed idea that natural laws, natural dispositions and
natural causes cannot possibly have anything to do with the
supernatural. That this fixed idea is unreasonable has been shown in what
we have said before. In the second place: the sudden, surprising and, for
just that reason, improbable character of the phenomenon of revelation as
it is presented to faith in authority. But to put the improbable on a
level with the impossible is unreasonable. As regards genius, for
example, no sharp boundary lines can be drawn between what, by organic
disposition, is owed to the activity of the brain, what is owed to
self-activity, and what is owed to inspiration. If it is so with the
inspiration of genius, how can one demand that in the miracle of
inspiration a critical boundary should let itself be drawn between what
is personally imparted by God's spirit and what lies beforehand in the
human soul? On the other hand, it is an inescapable demand of reason, as
regards miraculous as well as inspired genius, that the laws of nature
and of spirit, if in highly different ways, must find their fulfillment
in both.

\opage{40}This casts a new light on the miracle of revelation and its
appropriation. In antiquity the impression of the supernatural was
strong, but the development of reason comparatively slight; in our day
rational insight into the connection of natural things is rich and
clear, but the impression of the supernatural poor, weak and vanishing.
If it is to be possible for the consciousness of the present age to
bring about a psychological equilibrium between the religious and the
rational, it can happen only through an interaction between faith in
authority and personally free appropriation.

The word of revelation, with its wonders and holy facts, must, despite
the dissolutions of scientific criticism, remain standing for the
religious consciousness, unchanged, unchallenged. Reason must regard the
holy pictures of life as original paintings by old masters and not allow
the colors to be smudged and the features effaced by rationalist
fingers. Scientific authority these pictures do not have; but authority
of faith, authority of revelation, they have and will have until the end
of the world.

The authority of faith, too, must be dissolved; but, be it noted,
through a criticism of faith, into self-understanding and personal
appropriation. The way in which the dissolution takes place, the depth at
which the appropriation proceeds, differs in proportion to the
individuals' constitution of mind and the stage of formation at which
they stand.

As regards constitution of mind, there are souls of noble nature and
rich gifts whose religious need is so slight that they have no sense for
a personal spirit, but instead an original, living sense for the
impersonal universal, so ingenious a sympathy with the all-life in
nature that we must, in a special sense, call them \opage{41}nature-souls. The
English poet Shelley, Byron's friend, was such a nature-soul. In his
famous Ode to the West Wind he breathes out the need and longing of his
being. ``O, lift me as a wave, a leaf, a cloud! I fall upon the thorns of
life, I bleed. A weight of heavy hours has chained and bowed one who was
all too like you---untamable and swift-footed and proud. Make me your
lyre, as the forest is; what harm is it that my leaves fall like its
own! The might of your tumultuous harmonies will wrest from us both a
deep autumnal tone, sweet though sorrowful. Be you, O free spirit, my
spirit; be me, you violently storming one! Drive my thoughts, when I am
dead, over the universe like withered leaves, that they may call new
growth to life, and, conjured by this poem, scatter my words among
mankind like ashes and sparks from an unextinguished hearth!'' % Shelley, Ode to the West Wind, ll. 53--67 (the close of stanza 4 and stanza 5), given in Nielsen's Danish prose: translated from his Danish,
% which departs from Shelley at two points that matter here: "Be thou, Spirit fierce, / My spirit!" is "Vær Du, o frie Aand,
% min Aand", and "Drive my dead thoughts" is "Driv mine Tanker, naar jeg er død".
---Again, there are souls in whom the sense for inwardness is so strong,
the drive of self-preservation so ideal, the need for knowledge so deep,
that they find rest only when they have reached the stronghold of
all-love and feel themselves rooted in the eternal ground of the
self-thought. These spirited souls of light will not fly off with the
storm and be dissolved in the night of the elements; they will be
citizens of the kingdom of spirit and freedom. They are not Aeolian
harps that sound at the breath of the winds; they are instruments of
selfhood, tuned in joy and sorrow for holy harmonies, organs of will for
the self-conscious power that makes the storms its angels and the flames
of fire its ministers.

Which of the two kinds of soul, now, are the genuine, original human
souls? That cannot be settled in rational psychology; it is a faith that
must be tested in life.

If we take into consideration the various stages of development and
\opage{42}formation at which individuals within the circle of faith in revelation
can stand, the conception of what has been handed down will change with
each stage, and the religious consciousness will change color. The faith
of the child, the faith of the common people and the faith reflected in
reason point to different stages of development, and yet faith must be
the same in spirit and essence. For the educational system, and
especially for those who are appointed to be the teachers of youth, there
are serious difficulties to overcome here, psychological knots so
intricate that spiritual formation is needed to be able to untie them.

With the faith of the child it seems to go as of itself. The child's
soul is just as receptive to impressions from the supersensible as from
the sensible world. The wonderful vividness of the word of the Bible
appeals to its feeling and imagination, and the supernatural arouses no
offense: heaven is open, and God's angels ascend and descend upon the Son
of Man. But as childhood passes and the understanding develops, the
crisis sets in.---Let us, to make it vivid, picture to ourselves a
half-grown child and his teacher of religion reading together in the
Bible: of creation, Paradise and the Fall in the Old Testament, of
Christ's birth, his suffering, death and resurrection in the New. Here it
must of course be presupposed that the teacher believes sincerely in what
is being read about, and that the child is awake and well gifted. Until
the young pupil has reached his fourteenth year or thereabouts,
everything goes well. The boy takes pleasure in it and understands it all
so well, both the part about creation and the part about our first
parents, how beautiful it was in the Garden of Eden, and how grateful they
ought to have been, but that the serpent \opage{43}led them astray into
disobedience; it is almost as if he had lived through it himself. And now
the Gospel: the child in the manger, the angels' song of praise, the
shepherds' adoration! The childlike pupil understands it, loves it and
believes it as surely as if he had seen it with his own eyes. With the
earnestness of sorrow, in anguish of heart and grief, he follows the
suffering Savior over the brook Kidron to Gethsemane, to Caiaphas, to
Pilate's judgment hall; he lives through the crucifixion, he stands among
the pious women under the cross. The Savior dies; in the child's soul it
is as solemnly still as at the altar in a church. Then comes the morning
of the resurrection. The darkness vanishes, the guard flees, the angel,
clothed in light and lightning, rolls the stone away from the tomb.

With all this a change takes place little by little. The young pupil also
receives other instruction. As a pupil in a grammar school or a modern
school he acquires the first elements of science and makes good
progress, for he is gifted, has a ready grasp and a naturally sharp
understanding. He is now in his fifteenth year. The impressions from his
actual surroundings begin to gather themselves for him into a whole world
that is highly unlike the one he has received from the doctrine of
revelation; the thoughts of culture press in on him; he has secret
doubts and does not rightly know how. Meanwhile he is inwardly so taken
up with learning and remembering, and outwardly so immediately present
among his companions, that he has not yet had time to brood over the
connection; but in the religion lesson he has an objection at every
moment, which bobs up with an involuntary urge to ask. If the teaching
is conducted quite officially, he will take good care not to come out
with what is his own; the \opage{44}explanations he has been taught are more than
enough to put down all timid objections. But if he has confidence in his
teacher of religion and is so fortunately placed that he can talk with
him in all confidence, then misgivings, if in all unclarity, will steal
forth. Here is a turning point, where it is to be shown whether the
teacher is equal to his task.

If the teacher is a believer of the child's faith who, when it comes to
the point, knows no other way out than to appeal to the congregation and
the covenant of baptism and the disposition of the heart and the spirit
in the child's faith, then he may well, by the warmth of his words and
his heartfelt admonitions, make an impression on the young man, perhaps
even for a time seemingly lull his mind to rest; but the essence of
reason he does not see through, the principles of science he does not
grasp: how then is he to be able to help? A seed of doubt has fallen into
the young soul, and doubt shoots up and branches out, for doubt has
aerial roots like the mangrove, which forms the impenetrable forests in
tropical regions.

If the teacher is a high-church orthodox dogmatician, he will take the
matter in another way. ``Reason is blind to the holy; the heathenishly
puffed-up sciences of the present age are all on the way of error. God's
revealed word is to be believed literally. When it stands expressly in
Holy Scripture that God created the world in six days, then he created it
in six days, and what then do the objections of reasoning reason and the
teasings of the naturalists signify? God must surely know best himself
how long he took to create the world. God's word is infallible as God's
spirit is; had creation gone otherwise than it stands written, \opage{45}God's
spirit would surely have told us so plainly; for God is truthful, God's
holy spirit cannot lie. If heaven were not a fixed vault, how could it
stand in so many places that heaven opened? Where then would Christ have
gone when he went up to heaven? How could Paul be caught up into the
third heaven if there were not three heavens?''---But what is achieved by
such arguments other than that the young man now really begins to feel
himself a critic? The teacher can impress him with a forceful: ``Be quiet
and learn your lessons; you are still too green to dispute about such
things!'' The teacher can let him understand that he has a thorough
theologian before him, and give him hope that he himself may one day in
time become a thorough theologian and strike all the ``rudiments'' of
physics to the ground; but what does it help the poor pupil? The doubts
creep into hiding, but they sit in there.

If the teacher is a genial mixture of critic and theologian, half a
rationalist and half an orthodox, he will perhaps initiate the boy into
great secrets. He will regretfully concede to him that things are not
quite as right with Holy Scripture and our symbolic books as one might
wish. He will, with his hand on his heart, sincerely confess that there
are miracles in the New as well as in the Old Testament in which one
cannot believe at all: the miracle of the Gadarenes' swine, the miracle
of Balaam's ass, and many more. But withal he will then not neglect to
impress earnestly on his pupil that there really are also miracles in
which one ought unconditionally to believe---with qualification, that is.
The miracle of Pentecost, for example, is a \opage{46}principal miracle. There is
something grandly symbolic in it; that one must absolutely believe, so as
not to become a free thinker. To strengthen the young man at one and the
same time in criticism and in faith, the genial teacher will probably
also be so kind as to initiate him little by little into the ``critical
introduction'' and into ``apologetics.'' But what will be the
consequence? % the print has "hlive" for "blive" (PRINTED AS IS)
That the man-to-be, if with awakened sense he keeps his sharp
understanding unclipped, must end as a ``free thinker.'' If he will let
his understanding be circumcised, that is another matter; then he may
perhaps manage to become a theologian.

If the teacher is a man of spiritual formation, a personality who in the
fierce currents of the cultural consciousness has worked his way forward,
not to a sham child's faith with which one creeps back into the nursery,
but to a truly childlike faith, a faith purified in self-understanding, a
living faith in the revealed word, then he will know how to take the
matter as it is. He will neither rebuff the young man's doubt nor enter
into dispute about particulars; he will understand that it is once again
the two views of the world, that of antiquity and that of the present
age, that of revelation and that of reason, which now begin to wrestle
with each other in this soul. Like an experienced pilot he will help the
one still inexperienced over the first dangerous shoals, show him the
seamarks and orient him on the open sea. It is a psychological crisis he
has before him, and he will treat it with the superiority of the mature
man.---``What is the matter, my young friend, that you can no longer, as
before, believe in the wonder of creation, in Paradise and the Fall, but
are even beginning to doubt Christ's resurrection and ascension? A couple
of years ago we could both \opage{47}believe in them. I have not changed in this
respect. I still find the words of the Bible so telling, its expressions
so striking, its images so fresh, so full of life and power, that I am
rejuvenated by them and each time appropriate something new. It was the
same with you. I still see the tear that broke forth from you when, at
the place in the story of the Passion where Pilate says to the Jews:
behold \emph{what a man}! you gave vent to your heart and burst out: Had
He with the crown of thorns not come into the line of kings, how
detestable the long history of those with the crowns of gold would be! That
was a moment when you suddenly seemed to grow older than yourself in
being made young. But now you are beginning to grow old; what has
happened to you, let me hear!''

The young man takes courage and comes out with his confessions. One
day---he relates---we had physics in class; the teacher had grown warm
over the proofs of the law of gravity. It is, he said, a law so
necessarily grounded in the nature of bodies, so rational and eternal,
that no power either in heaven or on earth is able to suspend it. Then
one of my companions takes it into his head to ask whether Christ's
ascension did not make an exception. The teacher was silent, but with
such a \emph{look} and such an upturned glance that the whole class
laughed. I laughed too, although it was like a knife thrust in my heart.
This knife thrust gave my understanding a strange sharpness; it was as
if it all at once became angry and keen and wanted to cut my faith to
pieces. I now began to think over what I had read and heard about our
planetary system and our planet's insignificance in the universe, about
the animals of the primeval world and the millions, billions of years
that must have \opage{48}elapsed before man appeared on the earth. I am only a
beginner, and much of it is obscure to me; but I have an uncanny
presentiment, a secret suspicion, that Darwin's doctrine of man's
descent from some form of ape or other might well be the right one after
all. The apes have four hands, man only two; how probable is it not, then,
that the hind hands have been transformed into feet in the course of
time. Holy Scripture knows nothing of such things. The other day our
teacher of natural history drew our attention to the fact that no
serpent eats dust, although the opposite---he added with a sarcastic
smile---does indeed stand in the Bible. % the print has "sarkatisk" (PRINTED AS IS)
I sometimes feel as though I must either give up science or become a
free thinker; the first I cannot do, and the last I should be very loath
to do. Every time I come to think of it I feel very unhappy.

If you find yourself unhappy, the teacher answers, then the danger is
not so great after all. I was only afraid that you had become a
free-and-easy fellow and, in faith in the authority of opinion, had
already chosen the modern knowledge of the modern consciousness, as one
chooses a modern hat or a modern dress in a fashion boutique. You want
to pursue science, and in that you do right; for God has given you
abilities and gifts. The deeper you penetrate into the spirit of
knowledge and science, the greater respect you will have for such men as
Kepler and Newton, Laplace and Darwin; but you will also admire the
honesty with which such men always distinguish between what they know
and what they do not know. When you have come so far, you will bethink
yourself that a double limit is set to human knowledge: a \opage{49}relative one,
which can constantly be overstepped, and an absolute one, which cannot be
overstepped. If your teacher of physics once, at a public examination,
set you a problem from higher mathematics that had been solved by
Laplace, and demanded that you, without special preliminary study, solve
it with the help of your elementary school knowledge alone, would you
then regard this as a proof that the teacher himself is a great
mathematician, or would you not much rather regard it as an arbitrary
and unfair demand? Surely you will not suppose that the teacher who keeps
to the simpler problems, and in solving them uses a manner of
presentation that those not yet advanced can understand, always does so
because he himself knows no more; you will be modest enough to suppose
that he may well have a higher knowledge although he condescends to your
grasp. Now apply this to Holy Scripture, and ask yourself whether it can
be regarded as a proof of the divine spirit's ignorance in physics and
natural history % the print has "Psysik" for "Physik" (PRINTED AS IS)
that in its revelations about the origin of the world and the descent of
the human race it has condescended to the grasp of antiquity and spoken
in a language that men could understand. Your sound understanding will
tell you that what would be an unpardonable error in a textbook of
natural science can be a great merit in a religious communication, where
the natural is and remains always a symbol of the spiritual. You speak of
becoming a free thinker. If to free thought from ingrown prejudices and
preconceived opinions is the mark, then He who says: ``You shall know the
truth, and the truth shall make you free,'' is the greatest free thinker.
But if the game of letting one's \opage{50}thought go up in the air, as a boy lets
his kite go up in the air, and flutter in every wind, is to mean the
freeing of thought, then the free thinker is no more at home in science
than in religion. However loosely his thought flutters, it is
nevertheless not free; he has a string on it, as the boy on his kite: it
has secret orders to peck at what he does not like and to tear down what
goes against his instincts. In this sense, I hope, you will never become
a free thinker. If you love freedom---I mean freedom of the spirit---then
you cannot possibly suppose that true freedom should consist in making
spirit the slave of nature and laying it in the chains of necessity. The
God of revelation is the God of freedom; what do you think becomes of the
freedom of the man who denies the God of freedom and deifies necessity?
It looks grand that a man of free thought dares to go round in a ring
with the elements and free himself, almost as grand as that a dwarf dares
to creep into a giant's armor, and a Tom Thumb to fraternize with
Titans. If you love freedom, then the God of freedom is, after all,
personally your God, a God who is not outside you but in you as the
author of your being, and this God is spirit. Believe in the spirit of
freedom, of personality and of love, and the spirit will free you both
from unbelief's insult to the supernatural and from superstition's
contempt for the natural. Your doubt and unrest mean that you have
already set foot on the way of liberation; let me give you for the way
the words of Paul: ``When I was a child, I spoke as a child, I thought as
a child, I judged as a child; but when I became a man, I put away
childish things.''

We have dwelt so long on this transition because \opage{51}for our educational
system it is, at the present time, the most important of all. Could we
see the millions of souls who, although they received religious
impressions in childhood, have nevertheless since, amid the multiplicity
of experiences in the natural world, lost the original, the first simple
thing---whether, for the rest, they imagine that they have kept ``the
child's faith'' in the immediacy of authority, or have chosen
emancipation with ``the gospel of the flesh,'' or, in indifferentist
torpor, neither---could we see them in a psychological mirror, would it
not then appear that the incipient corruption stems from the period of
transition in which the child ripens into a youth and the youth into a
man, precisely the period when religious instruction is concluded? To
begin a religious instruction % the print has "Religionsunderviisnng" (PRINTED AS IS)
is comparatively an easy matter---any catechism parson and child-believing
infant-school Peter can do that; but to bring upbringing to such a
conclusion that a religious-personal self-knowledge, in union with a
truly rational cognition of the world, can be said to have been founded
is a difficult task, whose solution is conditioned by spiritual
formation.

\lecture{Fourth Lecture.}
\opage{52}It was in November 1842 that the pseudonym Victor Eremita's
\emph{Either/Or} struck into our literature like a bolt from a clear sky.
This stroke of lightning can truly be said to mark a turning point in the
development of the life of the spirit in Denmark. The very first lines of
the preface are characteristic. ``It may at times have occurred to you,
dear reader, to doubt a little the correctness of the familiar
philosophical proposition that the outer is the inner, the inner the
outer. You yourself have perhaps kept a secret which you felt, in its
joy or in its pain, was too dear to you for you to be able to initiate
others into it. Your life has perhaps brought you into contact with
people of whom you suspected that something of the sort was the case,
without your power or your allurement being able to bring the hidden into
the open. Perhaps none of these cases fits you and your life, and yet you
are not unacquainted with that doubt; now and then it has hovered past
your thought like a fleeting figure.''---There is something spectral in
this address that heightens expectation. The hermit with the fine, wise
face and the wonderfully clear, penetrating \opage{53}eyes seems to stand before
one large as life, with a riddle on his lips, a strange discovery, and to
open a secret hiding place. The voice is foreign and yet familiar, not
loud but attractive, sympathetic; it caresses the thought, it captures
attention, it whispers that there is something within the soul which is
called \emph{inwardness}.

The year after, another pseudonym appeared, Johannes de silentio, with
\emph{Fear and Trembling}. He foresees his fate in an age when one ``has
struck out passion in order to serve science.'' ``Love,'' he says, ``has
after all its priests in the poets, and at times one hears a voice that
knows how to keep it in honor; but of faith not a word is heard; who
speaks to the honor of this passion? Philosophy goes further. Theology
sits rouged at the window and courts its favor, offers its loveliness for
sale to philosophy.'' One notices from the pseudonym's mood that the
passion of faith has brought his mind to a lyrical boil. The year after,
a third pseudonym came forward, Johannes Climacus, with the philosophical
question: How far can the truth be learned? Here, he says, one meets the
difficulty to which Socrates in the Meno draws attention as ``a
contentious proposition'': that a man cannot possibly seek what he knows,
and just as impossibly seek what he does not know; for what he knows he
cannot seek, since he knows it, and what he does not know he cannot seek,
for he does not even know what he is to seek.---Everyone will understand
that by this pointing to Kierkegaard's pseudonyms the whole of
S. Kierkegaard's literary activity is meant, and that what we in this
connection lay weight upon, as regards spiritual formation, must be his
peculiar conception of faith and knowledge, religion and reason,
\opage{54}Christianity and civilization.

The last part of \emph{Either/Or} closes, as is well known, with ``the
upbuilding that lies in the thought that against God we are always in the
wrong.'' Our life brings us into a manifold relation to other people.
Some of these love right and justice; others do not, they do wrong. ``Your
soul is not hardened against the suffering they thus inflict on you, but
you search and test yourself, you assure yourself that you are in the
right, and you rest calmly strong in this conviction; however much they
offend me, you say, this peace they shall not be able to rob me of, that
I know I am in the right and that I suffer wrong.'' Our life brings us
into a manifold relation to other people; to some we are drawn with a
more inward love than to others. ``Now if such a person, who was the
object of your love, did you wrong, is it not true that it would pain
you, that you would carefully test everything, but that you would then
say: I know in myself that I am in the right; this thought shall calm me?
Oh, if you loved him, it would not calm you; you would search out
everything,\dots{} you would wish that you might be in the wrong, you would
seek to see whether you could not find something that could speak in his
defense, and if you did not find it, you would first find rest in the
thought that you were in the wrong.''

It is otherwise with our relation to God, for it is a relation of
infinity. Only in an infinite relation to God can doubt be set at rest;
only in an infinitely free relation to God can man's anxiety be
transformed into joy.
``In an infinite relation to God he is when he recognizes that God is
always in the right; in an \opage{55}infinitely free relation, when he recognizes
that he himself is always in the wrong. So then doubt is halted; for the
movement of doubt lay precisely in this, that at one moment he was to be
in the right, at the next in the wrong, to a certain degree in the right,
to a certain degree in the wrong, and this was to designate his relation
to God; but such a relation to God is no relation, and this was the
nourishment of doubt\dots{}. And should the thought that against God we are
always in the wrong not be inspiring, for what does it express but that
God's love is always greater than our love?'' It fills a person's soul
with strength to act, it makes him fervent in spirit; only when he
calculates finitely is the fire of the spirit put out. ``Even if you
yourself were the one who had to deny yourself your highest wish, you are
nonetheless glad; you do not say God is always in the right, for in that
there is no jubilation; you say: against God I am always in the wrong. If
what was your wish was what others and you yourself in a certain sense
must call your duty; if you had not only to deny your wish but in a way
to betray your duty; if you lost not only your joy but even your honor,
you are nonetheless glad; against God, you say, I am always in the wrong.
If you knocked but it was not opened; if you sought but did not find; if
you labored but received nothing; if you planted and watered but saw no
blessing; if heaven was shut and the testimony failed to come---you are
nonetheless glad in your work; if the punishment that the fathers'
brothers had called down came upon you, you are nonetheless glad, for
against God we are always in the wrong.''---This is the wonder of faith,
the tongue of inspiration, the proclamation of the ideal in the direction
of inwardness.

In ``Fear and Trembling'' Abraham is presented as the father of faith;
the point is the sacrifice of Isaac. The poetic gift \opage{56}with which the
presentation is endowed will always be an object of aesthetic admiration,
and the dialectical sharpness with which ``the absurd'' in faith is
driven to a higher madness will for a long time to come give matter for
religious reflections and ethical discussions. ``They rode in silence for
3 days; on the morning of the 4th day Abraham still said not a word, but
lifted up his eyes and saw Mount Moriah in the distance. He left the
young men behind and went on alone with Isaac by the hand up the
mountain. But Abraham said to himself: `I will not, after all, hide from
Isaac where this way is leading him.' He stood still, he laid his hand on
Isaac's head in blessing, and Isaac bowed to receive it. And Abraham's
countenance was fatherliness, his gaze was mild, his speech admonishing.
But Isaac could not understand him; his soul could not be lifted up; he
clasped Abraham's knees, he pleaded at his feet, he begged for his young
life, for his fair hope; he recalled the joy in Abraham's house, he
recalled the sorrow and the loneliness. Then Abraham raised the boy up and
walked with him by his hand, and his words were full of comfort and
admonition. But Isaac could not understand him. He climbed Mount Moriah,
but Isaac did not understand him. Then he turned away from him for a
moment, but when Isaac saw Abraham's face the second time, it was
changed; his gaze was wild, his form was terror. He seized Isaac by the
breast, threw him to the ground and said: `Stupid boy, do you think I am
your father? I am an idolater. Do you think it is God's command? No, it
is my desire.' Then Isaac trembled and cried out in his anguish: `God in
heaven, have mercy on me; God of Abraham, \opage{57}have mercy on me; if I have no
father on earth, then be you my father.' But Abraham said softly to
himself: `Lord in heaven, I thank you; it is better after all that he
believes I am a monster than that he should lose faith in you!'\,''---That
is one possibility of mood in ``Fear and Trembling.''

Yes---one will perhaps say---deep and dark are the depths of the sea, but
deeper and darker still are the abysses of the mind that open when faith
conjures in uncomprehended passion, and melancholy and eccentricity and
strange fancy keep one another company to accompany Abraham to Mount
Moriah and stare him in the eye just as he draws the knife. Abraham
stands as a venerable figure wrapped in the mists of antiquity; the great
features stand out clearly in the holy legend so long as he is seen at a
distance; but if one ventures to approach the place of sacrifice and
make his act present to oneself, then the understanding is lost, the
riddle of revelation, the terrible absurdity of faith, confounds its
thoughts. This misgiving is, however, as regards the essence of faith, at
least provisionally removed by the fact that the author has drawn for us
another exemplar, a figure as if taken out of the present age, a knight
of faith, at one and the same time hero and philistine. ``Here he is. The
acquaintance is made, I am introduced to him. The moment I first set eyes
on him, I instantly thrust him from me, myself jump back, clap my hands
together and say half aloud: `Good Lord! is this the man, is it really
he? he looks just like a tax collector.' Yet it is he after all. I draw a
little closer to him, watch the least movement, to see whether there
might not show itself \opage{58}a little heterogeneous fragment of telegraphy from
the infinite, a glance, a look, a gesture, a sadness, a smile, that
betrayed the infinite in its heterogeneity with the finite. No! I examine
his figure from top to toe, to see whether there might not be a crack
through which the infinite peeped out. No! He is solid through and
through. His footing? It is vigorous, belongs wholly to finitude; no
spruced-up townsman who goes out to Frederiksberg on a Sunday afternoon
treads the ground more firmly; he belongs wholly to the world; no
philistine can belong to it more. Nothing is to be discovered of that
foreign and distinguished nature by which one recognizes the knight of
infinity. He takes delight in everything, takes part in everything, and
every time one sees him take part in something particular, it is done
with a persistence that marks the earthly man whose soul clings to such
things. He attends to his work. When one sees him then, one would think
he was a clerk who had lost his soul in Italian double-entry
bookkeeping, so punctual is he. He takes a holiday on Sundays. He goes
to church. No secret glance or any sign of the incommensurable betrays
him; if one did not know him, it would be impossible to pick him out from
the rest of the crowd; for his sound, vigorous singing of the hymns
proves at most that he has a good chest. In the afternoon he goes out to
the woods. He delights in everything he sees, in the throng of people,
the new omnibuses, the Sound---when one meets him on the Strandvej, one
would think he was a shopkeeper's soul out on a spree, just so does he
enjoy himself; for he is not a poet, and I have in vain tried to spy out
in him the poetic incommensurability. \opage{59}Toward evening he goes home; his
gait is as untiring as a postman's. On the way he thinks that his wife
surely has a special little hot dish for him when he comes home, for
example a roast head of lamb with vegetables. If he met a kindred spirit,
he could go on all the way to the East Gate conversing with him about
this dish with a passion befitting a restaurateur. As it happens, he
does not own 4 shillings, and yet he firmly believes that his wife has
that delicious dish for him. If she has it, it would be an enviable sight
for distinguished people, and an inspiring one for the common man, to see
him eat; for his appetite is stronger than Esau's. His wife does not have
it---strangely enough---he is exactly the same. On the way he passes a
building site; he meets another man. They talk together for a moment; in
an instant he erects a building; he has all the powers needed for it at
his disposal. The stranger leaves him with the thought that he was surely
a capitalist, while my admired knight thinks: yes, if it came to that, I
could easily get it. He lies at an open window and looks at the square
on which he lives, at everything that goes on: that a rat slips under a
gutter plank, that the children play; everything occupies him with a
calm in existence as if he were a girl of 16. And yet he is not a genius;
for I have in vain tried to spy out the incommensurability of genius in
him. He smokes his pipe in the evening; when one sees him, one would
swear it was the grocer across the street vegetating in the twilight. He
lets things take their course % the print has "lade 5" for "lader 5" (PRINTED AS IS; Kierkegaard: "Han lader 5 være lige")
with an insouciance as if he were a frivolous good-for-nothing, and yet
every moment he lives he buys the \opage{60}opportune time at the dearest price,
for he does not do even the least thing except by virtue of the
absurd.''

The question of the absurd thus becomes, in the very highest sense, a
question of spirit, indeed a question of blessedness. In the
``Concluding Unscientific Postscript to the Philosophical Crumbs,''
Johannes Climacus presents the problem thus: ``The individual's eternal
blessedness is decided in time through the relation to something
historical, which is furthermore historical in such a way that there is
included in its composition that which by its essence cannot become
historical, and so must become so by virtue of the absurd.'' What this
means is shown vividly, from an exclusively Christian standpoint on an
evangelical-historical foundation, by Anti-Climacus: ``\emph{Practice in
Christianity}.'' The absurd, which to faith is a blessedness, is to
unbelief an offense. Christ, the object of faith, is the possibility of
offense: ``1) the possibility of offense that does not relate to Christ
(the God-man) but to him as a quite simply individual man who collides
with the established order. 2) The essential possibility of offense in
the direction of loftiness, that an individual man speaks or acts as if
he were God, says that he himself is God, and so in the direction of the
determination God in the composite God-man. 3) The essential possibility
of offense in the direction of lowliness, that he who gives himself out
to be God shows himself to be the lowly, poor, suffering, finally
powerless man.''

With this the standard that had hitherto held was annulled, and the
staff broken over the many learned, half-historical and half-philosophical
attempts to make Christianity comprehensible to human reason. \opage{61}Historical
investigations into the genuineness of the Bible, dogmatic speculations
about the Church's teaching and the like, may well furnish suitable
entertainment for the learned themselves, but for faith they mean
nothing. Faith is an offense to the understanding and foolishness to
reason; the believer believes the absurd: faith is a
paradox. How, then, are we to come to faith? By the way of inwardness
under the pressure of temporality, in the tribulations of finitude.
``Imagine, hidden in a simple setting, a secret compartment in which the
most precious treasure is deposited;---there is a spring that must be
pressed, but the spring is hidden, and the pressure must have a certain
force, so that an accidental pressure cannot be enough: so is the
disposition to faith, the hope of eternity, hidden in man's inner being,
and tribulation is the pressure. When the hidden spring is pressed, and
the pressure is strong enough, then the content opens in all its glory.
Or imagine a creeping animal that nevertheless has wings, which it can
use when it is driven to extremity, but for daily use it does not find it
worth the trouble to use them: so is the disposition to faith, the hope
of eternity, in man's inmost being: he has wings, but he must be brought
to extremity in order to discover them, or to get them and use them.''

By presenting faith as paradox, S. Kierkegaard gave the reasoning
understanding, with its self-important doubting of everything, a slap in
the mouth. With the paradox as presupposition he developed, in
``Works of Love,'' in ``The Sickness unto Death'' and in his many
upbuilding discourses, a complete cycle of religious categories, contour
lines of ideal circles, sound figures of a resounding eternity.---One
would have thought that a \opage{62}proclaimer of Christianity of such superior
gifts must have found support among all the teachers of Christianity,
the scholarly as well as the popular. But no! The theologians looked with
secret horror on the strange phenomenon, and there came over them a
feeling, an uncanny presentiment, that here was something with which
things did not hang together right. ``What, then? Surely there could
never be anything wrong with theology? Might it perhaps have to fear a
fire?'' Where are you going? Theology has nothing to fear; its clay is
fireproof, its stronghold is secure. The outhouses with sheep and cattle
and other livestock are, besides, insured. No, it is the poor author who,
with all his wit, his cleverness and his dialectic, is losing his mind.
That is a pity; otherwise one could have used him---not for church
history, for he lacks the learning and the grand overview for that; not
for dogmatics, for he lacks the speculative profundity and the sense for
the objective for that; but for apologetics: he could have won himself a
name as a Christian apologist. Now he must be given up.---And so the
watchword went out from headquarters that the whole extensive
Kierkegaardian literature was, until further notice, to be treated as a
heap of ``stray thoughts, aphorisms, notions and flashes.''

One would have expected that the clergy, the proclaimers of the word in
the congregation, would have made use of a helping hand by which they
could all at once be rid of the rationalist nuisance and have the
authority of the Bible, unclipped, for reliable use. But no! There must
be measure in all things. If the word of the Bible is to be really
upbuilding, it must be kept in a certain twilight of pastness. It does no
harm that there is a little uncertainty in the \opage{63}background, a little fear
of free thinking in people's minds; the proclaimer of the word then has
the advantage of being able to strengthen his hearers in the faith and to
speak an authoritative word of warning against ``the unbelief of these
times.'' For S. Kierkegaard, general reflections on ``those times'' and
``these times'' were idle dawdling. Anyone who will can become, in the
spiritual sense, contemporary with Christ in the same way as the
apostles. For about Christ, objectively seen, one gets to know nothing.
Christ is the object neither of pious admiration nor of criticizing
curiosity. Christ is the sign of contradiction, the possibility of
offense, faith's only-begotten. Christianity is not a doctrine; it is an
``existence-communication.'' Will you exist as a Christian or will you
not? That is the question. If you will not, you can go your way; if you
will---see, here is the pattern! So he held the New Testament so close
before the eyes of priest and congregation that with his
``contemporaneity'' he became really somewhat intrusive. He was answered
in the name of the Church that, as regards existence, there was nothing
amiss. ``We have the same Lord, the same faith, the same baptism, the
same Holy Spirit as the apostles. True, we cannot do such mighty works as
the apostles, nor live as economically as the apostles, for the times
have changed; but we rely securely on the Lord's word: Thou art Peter,
and upon this rock I will build my church, and the gates of hell shall
not prevail against it.''---Now the fortress was invested, and the
bombardment began.

``Imagine that there was a mighty spirit who had promised some people his
protection, but on the condition that they should present themselves at
a certain place \opage{64}to which it was dangerous to go---suppose now that these
people failed to present themselves at the appointed place, but went
about at home in their living room speaking inspired words with one
another about how the spirit had promised them his mighty protection, so
that no one should be able to harm them: is this not ludicrous?\dots{} That
the gates of hell shall not prevail against his church---these words of
Christ have lately been brought to mind again and again against me,
against my assertion that Christianity does not exist at all. My answer
is this. That promise does not help us in the very least; for the twaddle
we live in, as if it were being a Christian, is not at all what Christ
and the New Testament understand by being a Christian. Christendom is not
at all the church of Christ; nor do I say that the gates of hell have
prevailed against the church of Christ, by no means. No, I say that
`Christendom' is a twaddle that has clung fast to Christianity like a
cobweb on a fruit, and that is now so kind as to want to mistake itself
for Christianity.''

We strive and struggle, we yearn and seek. What is it we seek? ``First
the Kingdom of God.'' (A Sort of Short Story.)
Ludvig From, candidate in theology---he is seeking. And when one hears
that a ``theological'' candidate is seeking, one needs no lively
imagination to understand what it is he seeks: naturally the Kingdom of
God, which one is, after all, to seek first.---No, it is not that after
all; what he seeks is a royal living as a priest; and, as I shall
indicate with a few strokes, a very great deal has first happened before
he has got that far. First he went to grammar school, from which he was
then sent up to the university. Then he first took 2 \opage{65}examinations and,
after 4 years of study, first took his degree examination. He is then a
candidate in theology; and one might perhaps think that, after first
having got through all this, he can at last come to work for
Christianity. Yes, on the guild's terms. No, first he must spend
\textonehalf{} year at the seminary; and when that is over, there can be no
question of being able to apply during the first 8 years, which must thus
first be got through. And now we stand at the beginning of the story:
the 8 years have elapsed; he is applying. His life, which hitherto cannot
be said to have had any relation to the unconditioned, suddenly takes on
such a relation: he seeks unconditionally everything; he writes one sheet
of stamped paper full after another, runs from Herod to Pilate, commends
himself both to the minister and to the porter; in short, he is wholly in
the service of an unconditioned. Indeed, one of his acquaintances who has
not seen him for the last couple of years thinks, to his astonishment,
that he discovers that he has grown smaller, which can perhaps be
explained by its having gone with him as with Münchhausen's dog, which
was a greyhound but from all the running became a badger. So 3 years
pass. Our theological candidate really needs rest, needs, after so
enormously strenuous an activity, to be put out of activity, or to come
to rest in an office and be cared for a little by his future wife---for he
has meanwhile first become engaged. At last---as Pernille says to
Magdelone---the hour of his ``redemption'' strikes, so that with the whole
might of conviction he will be able, from his own experience, to
``witness'' to the congregation that in Christianity there is salvation
and redemption: he gets a living. What happens? On obtaining still more
precise \opage{66}information than he has hitherto had about the income of the
living, he discovers that it is about 150 rix-dollars less than he had
thought. Now all hell is loose. The unhappy man almost despairs. He has
already bought stamped paper to submit a petition to the minister that
he be permitted to be regarded as not called---and so to begin again from
the beginning: when one of his acquaintances gets him to give this up. So
it rests there; he keeps the living. He is ordained---and the Sunday has
come when he is to be presented to the congregation. The rural dean by
whom this is done is a more than ordinary man; he has not only (what most
priests have, and as a rule in a degree the greater the higher they rise
in rank) an unclouded eye for earthly profit, but also a speculative eye
on world history, something he does not keep to himself but lets the
congregation benefit from. He has, ingeniously, chosen as his text the
words of the apostle Peter: ``Behold, we have forsaken all things and
followed thee,'' and now explains to the congregation that precisely in
times like ours such men are needed as teachers, and in connection with
this he commends this young man, of whom the dean knows how near he came
to withdrawing for the sake of the 150 rix-dollars.---The young man now
himself mounts the pulpit, and the gospel of the day is (strangely
enough!): Seek \textbf{first} the Kingdom of God. He preaches his sermon. ``A very
good sermon,'' says the bishop, who was himself present, ``a very good
sermon; and it produced a real effect, the whole passage about `first' the
Kingdom of God, the way in which he emphasized this `first.'\,'' ``But
does Your Reverence think, then, that there was the desirable agreement
here between \opage{67}the speech and the life? On me it made an almost satirical
impression, this `first.'\,'' ``What an absurdity; he is, after all,
called to proclaim the doctrine, the sound, unadulterated doctrine of
seeking first the Kingdom of God; and that he did very well.'' % the quotation marks of this exchange (pp. 66--67) are incomplete in the print; set here as in Kierkegaard's text

By looking away from life and talking about the doctrine, the sound,
unadulterated doctrine, one obtains ``new-fashioned religious
safeguards.'' ``Once, in a long, long vanished time, people understood
the matter thus: it was demanded of him who would be a teacher of
Christianity that his life afford the guarantee for what he taught. One
has now long since got away from this; the world has become wiser, more
earnest, has learned to despise all this pettiness and morbidity about
the personal, has learned to covet the objective alone---now it is
demanded that the teacher's life afford the guarantee that what he says
is commerce, dramatic festivity, divertissement, purely objective. Some
examples. Is it this you will talk about, that Christianity, the
Christianity of the New Testament, has a predilection for the single
state---and you yourself are a single man:
my dear man, that is nothing for you to talk about; the congregation might
think it was earnest, become uneasy, or it might feel offended that you
so unbecomingly mix your personality in. No, it will be a long while yet
before you can become fit to talk about it in earnest in such a way that
you truly satisfy the congregation. Wait until you have once already got
your first wife in the ground and are well along with the second: then
the time has come for you; then you step forward before the
congregation, teach and ``witness'' that Christianity has a predilection
for the single state---and you will satisfy completely; for your \opage{68}life
affords the guarantee that it is nonsense, commerce, or that what you
say is: interesting.---Is it this Christian thing you will talk about,
that Christianity teaches contempt for titles and orders and all the
fooleries of honor---and you yourself are neither a person of rank nor
anything resembling it: my dear man, that is nothing for you to talk
about; the congregation might think it was earnest, or feel offended by
the lack of formation in thus forcing your personality upon it. No, wait
until you have first got yourself a heap of orders, the more the better;
wait until you yourself drag a string of titles about with you, so that
for the number of titles you hardly know yourself what your name is; then
the time has come, then you step forward, preach and ``witness''---and you
will undoubtedly satisfy the congregation; for your life affords the
guarantee that it is dramatic entertainment, an interesting forenoon
diversion. Is it this you will talk about, that to proclaim Christianity
in poverty is the true Christian proclamation---and you yourself are
literally a poor devil: % the print has "Diævel" for "Djævel"
my dear man, that is nothing for you to talk about; the congregation might
think it was earnest, become frightened, feel altogether out of humor and
be in the very highest degree uncomfortably affected by poverty coming
so close to one's life. No, first get yourself a fat living, and when you
have had it % the print has "har hav" for "har havt"
so long that you will soon advance to a still fatter one: then the
opportune time has come; then you step forward before the congregation,
preach and ``witness''---and you will satisfy completely; for your life
affords the guarantee that the whole thing % the print has "de Hele" for "det Hele"
comes to a jest, such as earnest men may wish for, in the theater or in
church, a refreshment for gathering new strength to---\opage{69}make money.---And
that is the way God is worshipped in the churches! And then these velvet
and silk orators weep; the voice fails, choked in tears!''

It is told in the Gospel that an angel at certain times went down into
the pool of Bethesda and troubled the waters, so that whoever among the
many sick first stepped into the bath thereafter was healed. Here too is
one who troubles the waters, the waters of faith, of thought and of
language; but is he an angel? He seems rather to be a demon; for anyone
at all among ``the healthy'' who ventures a leap into the troubled waters
must surely become lame in both soul and body. What, then, is he, what
does he want? ``What I want? Quite simply: I want honesty. I am not, as
people with good intentions have wanted to present me---for the view of
me taken by embitterment and rage and impotence and twaddle I cannot take
into account---I am not a Christian severity over against a given
Christian leniency. By no means; I am neither leniency nor severity---I am
a human honesty.''

The highest question of spirit and of culture has been set in motion.
Who is right, then, S. Kierkegaard or his opponents? The question goes
deep, and is moreover so sharpened and pointed that it sets its spear
point against the breast of the age and forces it to stand before an
Either/Or. One can turn one's eyes away from the question and act as if
it were forgotten or did not exist at all. But it exists all the same,
and will not be turned away. When one reads the collected newspaper
articles, Kierkegaard's attack and the clergy's defense, and takes
``the Moments'' along with them, there takes shape out of the
proceedings a peculiar picture, a situation that recalls \opage{70}the hero and
the chorus in ancient tragedy. This drama in our Danish literature will
surely pass into foreign literatures and be performed on the great stage
of the world. In a hundred years the lover of the paradox, with his works
of love and his ``1000 priests,'' will stride across the stage,
dramatically and with full music, like Don Juan with his 3000 mistresses.

\lecture{Fifth Lecture.}
\opage{71}\emph{The Christianity of the New Testament no longer exists}. That is
S. Kierkegaard's assertion, his ``only thesis.'' And the proof? It is
short and pithy. Christianity is ``existence-communication.'' Religious
existence in the present age and religious existence in the time of the
apostles are so diametrically opposed to each other that the one can be
regarded as a parody of the other. ``If what we understand by being a
Christian really is being a Christian, what then is God in heaven? He is
the most ludicrous being who has ever lived, his word the most ludicrous
book that has ever seen the light: to set heaven and earth in motion (as
he does in his word), to threaten with hell, with eternal
punishments---in order to attain what we understand by being Christians
(and we are, after all, true Christians): no, anything so ludicrous has
never occurred! Imagine that a man came up to a person with a loaded
pistol and said to him: I will shoot you down; or imagine, what is still
more terrible, that he said: I will seize your person and torture you to
death by the most horrible manner of death if you do not (now pay
attention, here it comes!)---if you do not make your life here on earth
\opage{72}as profitable and as full of enjoyment as is possible for you: then this
would surely be the most ludicrous speech\dots{}. The most horrible kind
of blasphemy is the one `Christendom' is guilty of: to transform the God
of spirit into---a ludicrous piece of twaddle; and the most spiritless
kind of worship of God, more spiritless than anything that is or was
found in paganism, more spiritless than worshipping a stone, an ox, an
insect as God, more spiritless than everything, is: to worship under the
name of God---a twaddle-head!'' % the print closes the quotation after „jeg skyder Dig ned“ and reopens it at „saa er dette“ (pp. 71--72); set here as one quotation

\emph{When did the Christianity of the New Testament cease to exist}? As
people gradually ceased to renounce the world, to sacrifice life for the
faith and to suffer for the doctrine. It still existed, then, in the time
of John Huss; Luther, too, must surely be counted among the Christians of
the New Testament? Yes, in a way; but if we take it strictly,
Christianity has not really come into the world at all. It stopped with
the pattern and at most the apostles; but these already proclaimed it so
strongly in the direction of propagation that here already the trouble
begins\dots{}. Only divine authority could so impress the human race that
it became an unconditioned earnest to will the eternal unconditionally.
Only the God-man would be able to endure, if one will think of it, to
work unconditionally through 1000 years and again 1000 years for the
propagation of the doctrine by proclaiming it, even if he did not get a
single disciple, if he could get them only by altering the condition. The
apostle, after all, has some selfish need of the relief of getting
followers, which the God-man does not have, who does not selfishly need
followers and therefore has only the price of eternity, no market
price.'' % the print has this closing quotation mark and no opening one

\emph{And whose, then, is the fault that the Christianity of the New
Testament no longer exists}? \opage{73}Men's fault, naturally. ``In the New
Testament, according to Christ's own teaching, being a Christian is,
humanly speaking, sheer torment, a torment compared with which all human
sufferings are almost only childish pranks. What Christ speaks of---for he
makes no secret of it---is crucifying the flesh, hating oneself,\dots{}
the most heartrending sufferings in hating father, mother, wife, one's
own child''\dots{} That was not at all to men's liking. So they resorted
``to hypocrisy, they falsified the definition of being a Christian. Being
a Christian, it was said, is sheer blessedness: `what were I, oh, what
were I, if I were not a Christian; oh, inestimable good to be a
Christian; yes, being a Christian is what first gives life real meaning,
its joys their savor and its sufferings their relief.'\,''

Here is S. Kierkegaard's \emph{manifest self-contradiction}.

First he defines Christianity as an existence-communication, a
supernatural communication of spirit, a gift of God's free grace, and
makes the expression more precise % the print has "præcicerer" (PRINTED AS IS)
to the effect that Christ, the God-man himself, the only true man of
spirit, is the only true Christian. The apostles, his nearest followers,
already betray, despite inspiration, on the very feast of Pentecost, a
dubious inclination to falsification, a ``selfish need of the relief of
getting many followers.'' And why? Because the spirit in the apostles was
not as strong as the spirit in Christ. So unconditionally is the
existence-communication dependent on the supernatural,
paradoxical-miraculous communication of spirit that the utmost exertion
of human power and will in this relation becomes a vanishing \opage{74}quantity. A
person can and should appropriate the divine in so far as it is
personally communicated to him, the single individual; but no person can
with all his zeal add one cubit to his stature, or appropriate even the
least of what is God's, if God himself will not communicate it to him.
On this presupposition, then, S. Kierkegaard with his ``human honesty''
should have come to the following result: the Christianity of the New
Testament no longer exists. Alas, how regrettable. But let us not lament
over it or make ourselves fruitless reproaches; for it is not men's fault
but God's inscrutable will. In the fullness of time God made an
experiment with a supernatural existence-communication, but let it drop
again shortly afterwards. How? For what reasons? Did his love grow cold,
did the spirit grow tired, did men make things so hard for him that he
grew weary of the whole thing? Ask not, mock not! It is, after all, what
it is meant to be---a paradox, an absurdity. With this explanation the
spokesman of the paradox could boldly have maintained his thesis and
consistently have finished with what he calls ``the twaddle of priests.''

But that was not the way he went. He took another run at it. He set forth
the pattern in order to enforce the demand for unconditional imitation.
Anyone who earnestly wills must be able to believe that the God-man has
existed. In faith in this existence the existence-communication, with its
supernatural conditions, is given once for all. Every believer who wills
can become contemporary with Christ, can step into his footprints and
follow after him. ``The relation is: the more you have to do with God,
and the more he loves you, the more, \opage{75}humanly speaking, will you become
unhappy for this life, the more will you come to suffer in this life.''
The supernatural, paradoxical-miraculous existence-communication is now
made dependent on man's mind and will to such a degree that there is not
even a shadow of change on God's side, for God is love and in the moment
has to do with everyone who will have to do with him. The thesis put
forward must then be defended thus: the Christianity of the New Testament
no longer exists; but it is not God's fault, it is men's fault. ``This I
know from experience in two ways. Partly from the fact that I cannot bear
the thought'' (the thought of love) ``and therefore merely, as one who
discovers, catch sight of this true Christian definition of being a
Christian, while for my own part I help myself to endure sufferings with
a far lighter thought, a Jewish one, not a Christian one in the highest
sense: that I suffer for my sins; partly from the fact that through the
circumstances of my own life I had to be led to become aware of
this''\dots{}. The ``human honesty'' is now this: the Christianity of the
New Testament no longer exists; for I, Søren Aaby Kierkegaard, have not
willed that it should exist. Had I willed to have more deeply to do with
the divine love and let it make me absolutely unhappy, wholly and
altogether unhappy in this life, then there would at least have been one
Christian: Christianity would then have existed in me. But as it is, I
have not willed it. Why not? Mock not, ask not about it; it is, after
all, what it is meant to be: the paradoxical, the absurd. My task I have
accomplished: to tear the mask off hypocrisy. The priests say that they
will what Christ wills, and to keep people \opage{76}in this illusion they put on
devout faces and falsify Christianity; I state truly what Christianity
is, and then I anoint my head and wash my face and say freely: Christianity
does not exist at all; I am not yet a Christian, because I do not yet
will to be one. That is my honesty toward the priests.

One sees that the last explanation does not fall short of the first in
effrontery and madness. Of such effrontery S. Kierkegaard has certainly
not made himself guilty. But by what has he avoided the extremity of
madness? By swerving from consistency, putting a button on the spear
point and confessing his frailty: I cannot bear the great thought of
Christ; the pattern is too great and I am too small; in other words: I am
still waiting for an existence-communication.

The Kierkegaardian relation to God is worlds apart from the relation to
God of the New Testament. The relation to God of the New Testament is an
immediate relation of revelation; the Kierkegaardian one is an infinitely
dialectical relation of reflection. The contradiction, the radical
self-contradiction, which S. Kierkegaard with all his reflecting never
brought up into consciousness, is this: that a derived relation of
reflection to God should be able to give the same existence as the
primitive relation of revelation. Certainly S. Kierkegaard thought and
acted, lived and suffered second to none in his age. His nearest posterity
has not yet properly bethought itself what work of the spirit he has
accomplished, and how infinitely much it owes him. But his reflection was
negative; the paradox, the absurd, was his only positive. The positive of
the paradox---that is its deep secret---is itself, in its inmost being,
negative. The blessed intoxication in \opage{77}suffering, renunciation, sacrifice
is a negative existence; God's love dissolves into sheer negative
determinations and is transformed at last into a soul-devouring depth of
the abyss that is more terrible than death. ``God is your mortal enemy;
he, love, he will out of love for you be loved by you: this means that
you must die, die away, otherwise you cannot love him. So there he sits,
omnipresent and omniscient as he is, and watches you, knows even the least
thing that goes on in you---that he knows, your mortal enemy! Beware of
wishing for anything, beware of fearing anything; for what you wish is
not fulfilled, but rather the opposite, and what you fear, the more you
fear it, the sooner it comes upon you. For he loves you, and he will be
loved by you; both out of love; but as soon as there is something you
wish, you are not thinking of him, just as when there is something you
fear; or if you bring him into connection with your wishing and fearing,
then you are, after all, not thinking of him in and for himself, that
is, you do not love him---and he will be loved, wills it out of
love.''---With such a love one lands not in heaven with Christ but in
Nirvana with Buddha.

As the revelation of truth that opens the holy springs of life,
Christianity can with reason be called an existence-communication, in so
far as it is an infinitely fruitful communication of life. But life and
existence are not, after all, synonymous words. The same life can express
itself in different existences: eternal life is the inexhaustible ground
of temporal existences. The eternal life in the essence of Christianity
does not stop at the stereotyped scheme of existence of a single age; it
\opage{78}enters, unchangeable as it is, into changes; eternal as it is, into time
and vicissitude; supernatural as it is, into natural conditions and
natural relations; blessed as it is, into earthly corruptibility, into
the unrest, sufferings and strife of worldly conflicts. Hereby the
relation between communication and appropriation is to be more closely
determined. Here there is a continued interaction between communication
and appropriation, but the factors of the interaction do not always
stand in the same relation.

At first the divine factor is almighty and overreaching. The Word becomes
flesh; the revealer of God steps forth as man, teaches and suffers, dies
and rises again in attestation of the personal victory of life over
death. He enters into glory; but just as he has taken his seat on the
throne at the right hand of the Father, the spirit comes down upon the
apostles as in a rushing mighty wind with tongues of fire, and the fire
kindles in these witnesses, and the tongues glow in the flames of
testimony. That was the great time of breakthrough; heaven was opened
and the graves were opened: the hidden wonder broke out, in this working
of the spirit, in manifest miracles.

Such is the primitive existence-communication, the Christian existence
in the New Testament. To measure the existence of the present age by
this standard, and to make ``Christendom'' answerable for having lost the
New Testament existence by its own fault, betrays an enormous
misunderstanding. The apostolic existence ceased when the period of
breakthrough was concluded.

There now follow periods in which the divine communication withdraws in
the form of the hidden \opage{79}wonder, so that the human factor of appropriation
may come forward in particular and arrive at conscious activity---the
Catholic and the old Protestant period.

If we look away from the differences in externals and fix our gaze on
the essence of appropriation alone, then the Catholic Church, by its
holiness of works, has undeniably done everything to exert the human
will, but it has also done everything to keep up the illusion of a
communication of spirit continuing in miraculous breakthroughs.---The
interaction between the divine and the human factor has thereby taken on
a stamp of mechanical finitude. The Protestant Church has pushed back the
notion of miraculous breakthroughs and reduced the communication of
spirit to a quiet, perennial wonder. It has not, indeed, wished to slacken
the power of appropriation, but it has nonetheless given the will a
one-sidedly inward-turning and self-dissolving direction in the dogma of
justification by faith.

Both churches are, each in its own way, in equally great embarrassment
with actual human nature and with existence as it works its way forward.
And why? Because, robbed of the gushing spring of the spirit that
supernaturally nourished the apostolic existence, they have nonetheless
held fast to the apostolic scheme of existence in the form of imitation.
We are, on the instructions of both, to renounce the world, mortify the
flesh, subdue the natural drives---in short, to regard the whole present
life as a misery that is good only for tormenting us up to blessedness in
the beyond. As regards humane existence, Catholicism, with its
thoroughgoing asceticism, its undisguised enmity toward free, rational
enlightenment and its monasticism's dark contempt for the natural ends
of life, is, however, more consistent \opage{80}than Protestantism. The Protestant
Church, while retaining its apostolic scheme, has wanted to adopt
humanity and take the beauty of culture into its embrace; that is an
ambiguity for which it has had to pay dearly. Where now is the New
Testament existence, where is the mortification of the flesh, where is
the Kingdom of God? ``Ludvig From, candidate in theology, he is
seeking.''---S. Kierkegaard's satire has a tangible effect; every stroke
of his scourge leaves bloody weals, precisely because he makes valid the
New Testament scheme of existence, solemnly acknowledged by his opponents
themselves.

The existence of the present age has peculiar demands which it neither
can nor will give up; it asserts the significance of natural life from
the ground up and demands that reason be acknowledged. If Christianity
really were a paradox, and faith an absurdity, then the Christianity of
the New Testament would certainly no longer exist; but, be it noted, for
the reason that it has never existed except in men's imagination, an
imagination that can hold out only until the paradox has been discovered.

What is the paradox? The paradox, with the absurd, comes to light when
two concepts that, in full transparency, show themselves to be mutually
contradictory are put together in unity. A square circle is an
absurdity. If one now says: in the highest regions of spirit and truth
all circles are square, and this is to be believed against all
understanding, then reason protests; we then have the paradox. If one
says: in the existence of the God-man something has become historical
that by its essence cannot possibly be historical, and so must have
become so by virtue of the absurd, and this is to be \opage{81}believed against all
understanding, then reason protests: we have the paradox. That the
contradiction in the square circle and the alleged contradiction in the
essence of the God-man are, however, absolutely heterogeneous
contradictions, anyone who will use his reason can see. The concept of
the circle and the concept of the square can be made clear in exact
determinations and corresponding forms of intuition. But the content of
the concept of God we cannot exhaust in exact determinations, nor that of
the concept of man either. When, then, faith holds fast the unity of both
concepts in the essence of the God-man, reason grasps that what is
thereby presupposed is not an absurdity but a mystery.

Perhaps ``From, candidate in theology,'' will remark that this is
just what he has long known from theology, and for just that reason is
seeking\dots{}. Let him seek! The mystery is not an old theological night
in which all cats are grey; it is not the indeterminable twilight in
which religion is stirred together with science, faith with reason, into
edifying explanations for unclear heads. The mystery is a limiting
concept in which faith and knowledge meet and part. It must be precisely
determined, by careful criticism of knowledge, how deep reason can reach
in the cognition of the natural; the absolute limit at which cognition
dissolves is the mystery. It must be precisely determined, by careful
criticism of faith, how deep the religious consciousness can go in the
cognition of the supernatural without coming into contradiction with
reason and thereby into conflict with itself; the absolute limit at
which it is dissolved is the mystery.

\opage{82}The miracle of revelation is a mystery, the essence of the God-man a
mystery. If one now fixes one's attention one-sidedly on the miraculous
outbreaks and the corresponding crises of existence, while arbitrarily
avoiding thinking of the intermediate determinations hidden in the
mystery, one does quite rightly get out of it that the whole is an
absurdity. If, with S. Kierkegaard, one applies such a scheme of
existence as the standard to the existence of the present age and demands
commensurability, one does quite rightly get out of it that faith and
Christianity, the acknowledgment of the absurd, the paradoxical, no
longer exist. But the mystery of revelation still exists, and will go on
existing.

``The ideals must be proclaimed!'' Right. But the question is how they
are to be realized and take shape as existence. Abraham still stands as
an exemplar of faith; but not directly by his scheme of existence. ``God
tempted Abraham and said''\dots{}. God, then, spoke to Abraham as a master
can speak to his servant. This scheme we cannot appropriate. At our stage
of existence, anyone who appeals to divine commands by immediate
prompting will be regarded as of unsound mind. On the other hand, that a
person, after conscientious self-examination and mature deliberation,
finds himself on religious grounds obliged before God, absolutely
obliged, to sacrifice even what is dearest and most precious to him in
the world, cannot be regarded as unsoundness of mind. ``God said to
Abraham: take Isaac, thine only son, whom thou lovest, and go\dots{} and
offer him for a burnt offering''\dots{} If this scheme is to hold
literally for the present age, then the God of believers becomes a
Moloch, and the highest existential expression of faith and obedience
\opage{83}would then be: to sacrifice one's children to Moloch. On the other hand,
a personal demand of the spirit, even if in its inner unconditionedness
it must come into strong collision with the notions of the age, may well
be conceived to be grounded and fully justified. He who, in fidelity to
the will of the spirit, with God alone as his confidant, is to satisfy
the demand, will learn what obedience is, what it is to believe and to
bear a responsibility. When he then, under the heavy pressure of the
heavy responsibility, lifts up his head and sees Abraham standing on
Mount Moriah with the sacrificial knife in his hand, then the ideal is
drawn on the dark ground of the mystery; he sees a shining, uplifting
example of faith and obedience.

``The ideals must be proclaimed!'' Yes. But surely not by their being
reduced to absurdity. % the print has "redigeres in absurdum"
The humorous figure that S. Kierkegaard's pseudonym has formed by
fusing knight of faith and philistine is, as drawn, sound to the core;
but by the mechanics of the absurd it moves on screws. ``He (the knight
of faith) drains in the infinite resignation the deep sadness of
existence, he knows the blessedness of infinity, he has felt the pain of
renouncing everything, the dearest one has in the world, and yet
finitude tastes to him just as good as to one who never knew anything
higher, for his remaining in finitude has no trace of a cowed, anxious
training, and yet he has this security in delighting in it as if it were
the most certain thing of all. And yet, yet the whole earthly figure he
presents is a new creation by virtue of the absurd. He resigned
everything infinitely, and then he grasped everything again by virtue of
the absurd. He constantly makes the movement of infinity, but he \opage{84}does it
with such correctness and assurance that he constantly gets finitude out
of it, and there is not a second in which one suspects anything else.''

This marvel is now illuminated in particular by means of a young swain
who puts the whole content of his life into his love for a princess,
although the relation is such that this love ``cannot possibly be
realized, cannot possibly be translated from ideality into reality.'' To
state exactly where the absurd lies, the knight of infinite resignation
is set over against the knight of faith. When the knight of infinity has
seen the impossibility of getting the princess, he resigns the fulfillment
of his wish. His love takes on a religious character and is transfigured
into a love for the eternal being, which indeed denies him the
fulfillment but nonetheless reconciles him in the eternal consciousness
that he possesses the beloved in a recollection that no actuality can
take from him. He makes the impossible possible by expressing it
spiritually, and he expresses it spiritually by renouncing it. The wish,
which would have led him out into actuality but was wrecked on the
impossibility, is now bent inward, but is not on that account lost,
nor forgotten either. From the point where % the print has "kvor" for "hvor"
the knight of infinity stops, the knight of faith makes one movement
more, more wonderful than everything, for he says: I believe nevertheless
that I shall get her, namely by virtue of the absurd, by virtue of the
fact that for God all things are possible. Thus the knight of faith is
``the only happy one, the heir apparent to the finite,'' while the knight
of resignation is ``a stranger and a foreigner.''

The cardinal point about which the movements turn is the boundary point
between the possible and the impossible. \opage{85}Each of the two lovers is agreed
with himself that it is, humanly speaking, impossible for him---not merely
improbable, in the very highest degree improbable, but absolutely
impossible---to get the princess. Let us now look more closely at what
this implies. For human understanding to be perfectly certain of the
impossibility, it must have an unmistakable sign. What sort of sign would
that be? That the royal father, even in full accord with his council of
state and the whole diplomatic corps, European, Asian, American and
Australian together, set himself against the match would not be the
unmistakable sign. The king might, after all, die, and the diplomatic
positions change. Even if the princess's heart were cold to her lover's
pleas, that still would not be the unmistakable sign. It would not, after
all, be impossible that her feelings might change; there are examples of
young girls who have given up their coyness and at last let themselves be
moved. But---let us posit this case---if the princess died: yes, then we
have the sign, for it is certainly, humanly speaking, impossible to marry
her when she is dead and buried. What does the knight of resignation do
now? He resigns infinitely the fulfillment of his wish, but keeps his
love in the pain of recollection. But what does the knight of faith do?
He gets the---dead and buried---princess, and lives with her in full
finitude, in external actuality, in joy and glory, by virtue of the
absurd. If this is to be understood literally, then one must undeniably
grant the pseudonym is right when he exclaims: ``the last astonishes me;
my brain turns round in my head at the thought of it. The movement of
resignation I make by myself, and what I gain by it is my eternal
consciousness \opage{86}in blessed understanding with my love for God. But by
faith I renounce nothing; on the contrary, by faith I get everything,
precisely in the sense in which it is said that he who has faith as a
grain of mustard seed % the print has "Sennopskorn" for "Sennepskorn"
can move mountains.''

Either reason or faith! Where faith begins, reason ceases; i.e., so long
as a person has still kept a little remnant of sound reason, infinite
resignation is his highest; he reaches only the point where faith should
begin. The believer has put the whole stretch of reason behind him and
goes further. For him the boundary between probable and improbable,
possible and impossible, no longer exists, for with God, after all, all
things are possible. So: either reason, and then without any faith, or
faith, and then without any reason; this is the Kierkegaardian dilemma.
But if we look more closely, might it not then appear that faith and
reason, although in the human sense two absolutely heterogeneous factors
of cognition, nevertheless have their ground of unity in the depth of the
divine being, in the mystery?

The pseudonym in ``Fear and Trembling'' assures us that he can make the
movement of infinite resignation by himself---without faith---in
``blessed understanding'' with his love for God. But a blessed
understanding with God is impossible when the thought that for God all
things are possible is excluded. It is said in ``The Sickness unto
Death'': God is this, that everything is possible. By denying this
thought one abolishes the concept of God and denies God; but a blessed
understanding in love for a denied God is an absurdity. If this absurd is
absolutely to be the mark of faith, then the \opage{87}knight of infinite
resignation already finds himself, in his way, within the domain of
faith. The absurd in the knight of faith is to be this: that in the midst
of finitude he knows how to come to terms with the difference between
what, in the human sense, is possible and what is impossible, between the
probable and the improbable, although he nevertheless holds fast to the
thought that for God all things are possible. The knight of faith is
therefore, in his way, a quite reasonable man. But if the knight of
resignation, precisely by being reasonable, is a believer, and the knight
of faith, precisely by being a believer, is reasonable, then the boundary
is, after all, removed.

Has the opposition in principle between faith and reason vanished, then?
By no means! Here the start is made from two absolutely heterogeneous
points of departure: the connection of things by natural necessity, and
the omnipotence of the divine will; but as one descends into the depth
from the one side and ascends out of the depth from the other, it appears
that the last and outermost boundary between possible and impossible
vanishes in the mystery. The ideal of reason is to bring all the
possibilities of actuality under the conditioned necessity of the
existence of all things; the ideal of faith is to bring necessity, with
the infinite chain of conditions, under the unconditioned freedom of the
divine will; but no mortal shall be able to point out the last and
outermost boundary where conditioned necessity ceases and absolutely
unconditioned freedom begins. The original, actual unity of freedom and
necessity is hidden in the mystery.

In presenting the ideals of faith the Kierkegaardian pseudonyms speak
with different tongues and contradict one another most forcefully,
precisely by virtue of \opage{88}the absurd. In Johannes de silentio the absurdity
of faith makes the believer the patriarch of the promises, the heir of
the blessings, the heir apparent to the finite, the most blissful man on
earth: Abraham gets Isaac by virtue of the absurd; the knight of faith
gets the princess---the princess dead and buried for the sake of
impossibility---he really gets her; they live together in joy and
gladness, by virtue of the absurd. But in Johannes Climacus and in
Anti-Climacus it goes otherwise. In Climacus, faith causes the reflecting
humorist, who would so gladly understand it and enter into it, the most
dreadful troubles of existence, by virtue of the absurd; and in
Anti-Climacus, the religious character who is and lives in the midst of
faith, faith makes man, humanly speaking, the unhappiest of creatures,
precisely by virtue of the absurd. But when the absurd is just as potent
to make a person happy as to make him unhappy, how can the difference
between the happy and the unhappy become the gauge of the strength of
faith?

The existence-scheme of sufferings, pains and renunciations can in
principle be just as misleading as that of joys and enjoyments. Christ
did not come into the world to intensify sufferings or to heighten joys;
that he nonetheless does both means that he came in the absolute sense,
for the sake of life and light and truth, for the revelation of God's
secret counsel, from the holy depth of freedom, the mystery of
righteousness and love.

It is quite another thing when the educative element is spoken of. From
the standpoint of upbringing and spiritual development, S. Kierkegaard
has \opage{89}unmistakably grasped what is right, in so far as, in his proclamation
of the ideals, he has again and again, with the penetrating sharpness of
a psychological observer and the convincing power of a painful
experience of himself, enjoined voluntary obedience under the law of the
spirit and preached ``the gospel of sufferings.''
But the existence of the present age and its spiritual conditions he did
not understand. He thought to overawe the established order by bringing
forward the New Testament demands of existence and, with colors of
glowing spirit, burning every feature of the scheme of the great
breakthrough into the consciousness of his contemporaries; and he did not
notice that his age, however run down it was by its ``new-fashioned
religious safeguards,'' was nevertheless right against him on one very
essential point. He thought to bring about a spiritual revolution, a
decisive catastrophe. Even on his deathbed he is said to have said: ``I
fall like a bomb, the bomb bursts, and then it catches fire.'' But the
commotion passed, and the catastrophe failed to come; the bomb burst, but
nothing caught fire. On the other hand, he undeniably succeeded, in a
quite striking and conspicuous manner, in surprising and catching the
whole official church establishment in the consequences of its own
presuppositions.

The single individual is not to have the right to mock the congregation
and put his subjective reflection in the place of the communication of
spirit; but let it only be said straight out: there is in the official
proclamation of Christianity something that must be changed from the
ground up. It is not the revealed word that is to be changed; it is the
appropriation, the free, personal appropriation of the rich truth, that
must be presented in a new light. If the thought of Christ is to take
root in souls and become the \opage{90}deepest, most powerfully moving thought of
life in the age of culture, then the images of the child's faith and the
notions of the faith of the common people will not suffice. It is a real
understanding of the compatibility of the religious with the rational on
which everything above all depends. The transformation of existence is
not catastrophic; it takes place---given an exertion of will---from
within, slowly and quietly, but surely and soundly, through upbringing and
progressing spiritual formation.

\lecture{Sixth Lecture.}
\opage{91}The existence of the present age has its own character; in bearing, gait
and physiognomy it distinguishes itself clearly and definitely from all
its predecessors. Its way of seeing is realistic; it keeps to the actual
as it stands and goes. Yet it by no means intends thereby to deny spirit
or to despise ideal demands; but it has a suspicion of all abstract,
high-flown ideals, whether they put on religious, ethical, political or
aesthetic forms. Ideals in the blue may be as lovely as they like; they
are nonetheless barren clouds, airy vegetations that cannot thrive in the
soil of actuality; they work harm when, in the name of eternity, they suck
the blood of temporality and drain the marrow from the present life
instead of giving it strength. The existence of the present age is by no
means without conscience, nor does it exactly suffer from scruples: its
conscience is like a fetus that has turned in the womb.

That human nature, on account of its doubleness (dualism) as sensible and
spiritual, constantly gives the impression that man lives as a frontier
soldier between two realms, a natural and a supernatural, is a fact. The
opposition between the spiritual and the \opage{92}bodily, the heavenly and the
earthly, will not let itself be effaced, however much one strives, for
the sake of harmony and unity, to weaken and suppress the one member. The
naturalistic endeavor, to glorify nature at the expense of spirit, and
the ascetic one, to help spirit to victory by killing nature, are
corresponding extremes. Neither of these extremes is able to beguile the
present age. The consciousness of the present age has criticism enough to
see that a naturalistic morality is not only just as hollow as
asceticism, but, in all its parading of the glory of sensuous life, still
more contemptible.

The ascetic, in the midst of his error, nevertheless has an unconditional
acknowledgment of spirit as the highest. His conception of the essence of
spirit may be poor and unclear; but once he has recognized that spirit is
will, he does not spare himself, but summons up everything to obey the
command of the spirit and satisfy its demands; he mortifies his flesh, he
stifles the natural lusts, he hardens his will by refining it in the fire
of the pains of renunciation. One may deplore his one-sidedness; but one
cannot deny him a certain esteem.

The naturalist, with ``the autonomy of the drives'' and ``the gospel of
the flesh,'' seeks his point of support in ``the modern consciousness.''
He can with a certain right appeal to the fact that natural life is the
foundation of the life of the spirit, but he distorts the truth by
reducing the alleged life of the spirit to a formal reflected gleam of
the consciousness of the essence and working of natural, finite,
corruptible life. He can make a good deal of fuss over the works of the
spirit that have already been accomplished and are yet to be accomplished
in the future before the ideal comes \opage{93}into its own and the work of freedom
is completed. But if we look more closely, the whole thing comes down to
a denial of the absolutely self-valid will of spirit, and so of spirit as
spirit. The divided forces of will are transformed into elementary forces
of nature. The great drama that the conflicts of these natural forces,
necessary by reason, produce receives its spirited stamp by being
beheld, criticized and enjoyed by ``the modern consciousness.'' The highest
task of mankind, the depth in all its longings, the reward of all its
efforts, the goal of its strife and striving, must then be placed in
raising itself up to the standpoint of ``the modern consciousness.'' What
the Church was for the believers of the Middle Ages, ``the modern
consciousness'' is for the knowers of the present age. The Church gave
indulgence for all manner of sins and even winked at gross excesses,
provided only that one condition was present: faith in the Church and in
the infallibility of its authority. If this condition was not present,
then woe to the offender; even if his life had the purity of the angels
and all the good works of a saint, he was nonetheless a heretic,
excommunicated and condemned to eternal perdition. The critical majesty
of ``the modern consciousness'' has stepped into the place of the Church.
Thought is free---that is to say, when it submits unconditionally to the
infallibility of ``the modern consciousness''; for this infallibility is
not merely an article of knowledge, it is an article of faith. The
devotees of science may have conflicting views, doubt reason, deify
reason, and make themselves guilty even of gross inconsistencies: no
matter! Criticism gives indulgence for all manner of contradictions and
winks at the frailty of thoughts, even at \emph{gross inconsistencies},
yet, be it noted, on one condition: \opage{94}that their authors sincerely profess
``the modern consciousness.'' If this condition is lacking, the guilty one
may pack up and hang himself on the tree of knowledge. The heretic is
condemned---not to stake and fire, for ``the modern consciousness'' holds
to what is humane---but to be expelled from science, robbed of the
future and set in the pillory of literature, for a terror and a warning
to others.

The consciousness of the present age, seen in its essence, by no means
coincides with ``the modern consciousness.'' The knowers of the present
age know very well how to distinguish between the passing and the
abiding. Truth is neither old-fashioned nor modern; it is eternal. The
eternal truth can change form and expression as the times change, but in
essence and ground it always remains the same.

The spirit of personal truth that reveals itself in religion can, at the
stage of the consciousness of the present age, be recognized in two
mutually corresponding forms of activity: the original form of
communication and the derived form of appropriation, the first
characteristic of the existence of breakthrough, the second decisive for
the existence of the present age. Never before in history has the
opposition between the two forms of existence stood out so clearly. The
relation has hitherto been more or less obscured by the fact that one has
partly wished to arrange the derived in the likeness of the original,
partly to give the derived the appearance of being the original itself
in its further continuation, and partly, finally, to remold the original
by accommodating it to the derived. Catholicism furnishes telling
examples of attempts in the first two directions. Apostolic holiness,
apostolic miracles repeat themselves in deceptive imitations; even the
Savior's life and holy wounds are \opage{95}copied in medieval color prints. The
old Protestantism let the original stand as the perfect, the
unattainable, and contented itself with appropriating the doctrine,
without being able to come to clarity about the transformation of
existence. While in actuality it entered upon an accommodation to the
changed conditions of the times and gave full acknowledgment to the
authority of the state and the manifold demands of civil society, it
nevertheless felt itself obliged in conscience to enjoin the literal
observance of the apostolic precepts, without daring to admit to itself
that the latter stood in glaring contradiction to the former. Only in the
period of the Enlightenment did people begin to take courage. By the
light of reason they saw how unreasonable it was to live directly by the
apostles' instructions. But what did they do then? They made the apostles
into ``reasonable'' men, gave their words a ``reasonable'' meaning, and
drew a plain bourgeois morality out of their teaching. The alleged
``perfectibility'' of Christianity had the consequence that the derived
was put in the place of the original: the copy became the original, and
the original the copy. The rationalist blunder was, however, too obvious
to hold its ground in the long run. The apostolic existence has stepped
forth again in all its peculiarity, and the question is how the existence
of the present age is to arrange itself so that the derived may, in spirit
and in truth, correspond to the original.

The first thing needful is that the opposition come out into full light.
The existence of revelation, with its ideal marks of character, remains
standing unchallenged. Over against this authority the existence of the
present age takes its stand on its historically determined foundation,
with clear consciousness \opage{96}of its peculiar spiritual conditions. The
existence of the present age knows in itself that it neither can nor will
set itself up to be a direct imitation of that ideal existence of
breakthrough. The ideal existence commends, for example, the single
state; but the existence of the present age maintains that marriage and
family life can, in the name of the spirit, be just as justified as
celibacy. When the ideal existence prefers poverty to wealth, the
existence of the present age is of the opinion that one can prefer wealth
to poverty without offending the ideal. When the ideal existence of
breakthrough looks down on political conditions of society as on
something indifferent in itself, and regards being a slave as just as
favorable a condition as being free, then the existence of the present
age declares straight out that the political condition is by no means
indifferent, that there must be fighting, in the power of the spirit, for
civil order and civil freedom, and that esteem and honor are better than
lowliness and dishonor.

What the present age thus declares is the sincere conviction of its
heart. By showing itself as it is and speaking its mind straight out, the
existence of the present age places itself in a truer relation to the
holy authority than by flattering and dissembling. Nothing can be more
odious to the spirit of truth than hypocrisy; the divine majesty is not
accessible to devout compliments.

It might indeed seem as though the existence of the present age, by
declaring itself in a direction so opposed to the ideal existence, meant
to defy authority and withdraw from the demands of the spirit. That is a
great misunderstanding. The meaning is that the two forms of existence,
the ideal existence with its original and the real existence with its
derived, must \opage{97}come out into full opposition, so that their unity may be
recognized in its depth and the bond of the spirit between them be tied
all the more firmly. What separates the existence of the present age from
the first original is something more than ``the unbelief of the world''
and ``the ungodliness of these times''; it is something over which the
human will has no arbitrary command. That breakthrough of the spirit did
not, after all, take place on account of men's holiness; rather, the
spirit broke through in order to sanctify men. The apostles were, by
their own confession, sinners, and the members of the apostolic
congregation were great sinners.

It is on the communication of spirit and the corresponding historical
conditions that the opposition essentially rests. In the existence of
breakthrough the self-communication of the divine spirit was so mightily
overflowing that it drove its chosen ones forward to voluntary testimony,
strife and struggle for the kingdom of light and truth. Over against this
kingdom of light and truth stood the kingdom of the world, the kingdom of
darkness and falsehood, which had to be fought with spiritual weapons.
What was hostile to spirit was external; suffering and tribulation, and
indeed discord within man himself, were conditioned in particular by this
externality: the spirit was willing, though the flesh was weak. In the
existence of the present age the communication of spirit is weaker, not a
miracle manifest in visible outbreaks but a quiet, perennial wonder. The
abrupt, catastrophic transitions from darkness to light, from unbelief to
faith, that characterize the period of breakthrough have been transformed
into gradual, step-by-step transformations, conditioned by the strength of
the will in the labor of appropriation. In the time of the apostles it was
the spirit that in a supernatural manner seized human souls, filled them,
lifted them, raised them above time and corruptibility. \opage{98}In the present
age, conversely, it is men who must seek the spirit present in its means
of salvation, seize it, hold it fast, and by inward appropriation work
their way up through the entanglements of temporality to truth and
eternal life. The swift movement from above downward has been converted
into a slow movement from below upward.

The kingdom of the spirit, which in the language of revelation is called
the Kingdom of God, is called in the language of the present age the
kingdom of humanity. Now if it could be shown that the Kingdom of God was
too narrow to embrace everything that belongs to the truly human in its
full extent, or if it could be established that the kingdom of humanity
had a foundation of its own on which it could be built, fortified and
extended without the help of the Kingdom of God, then the humane and the
religious would certainly fall apart. If the humane really fell outside
the religious, then it would be quite in order that the two realms of the
spirit should go on fighting each other to the death, for the complete
victory of the one would necessarily be the downfall of the other. But it
appears, on the contrary, that both realms, originally seen, have the
same foundation, the same highest end, so that no one can enjoy full
citizenship in the one without being a citizen of the other; here there
are not two realms but two forms of existence of one and the same
divine-human kingdom of the spirit.

When it nevertheless looks as though there really were two realms here
fighting each other, this cannot be ascribed to the constituting
fundamental law; it must be the fault of the citizens, who, by clinging
one-sidedly either to the one or to the other of the two forms of
existence, fall into discord with one another, form \opage{99}parties and plague the
realm with internal disturbances. Let us try to illustrate this by some
examples.

\emph{The laws of the spirit}. By taking his stand exclusively on the side
of the Kingdom of God and speaking the language of revelation, the
preacher of repentance gains the authority to assert the unconditional
validity of the holy laws. They are, after all, not commandments of men
but commandments of God. God himself has expressly given these laws; the
description of the divine legislation makes an imposing effect. Jehovah
thunders from Sinai, and the mountain trembles, and the consciences of
the hearers tremble: these are the commandments of the Almighty; woe to
him who dares to transgress them!---The history of the Pietists and the
Methodists furnishes telling examples of what stirrings of the spirit the
preaching of repentance is still able to call forth in recent times. A
holy authority is needed to awaken the natural man to earnest
self-examination. There is something awakening in severity itself: the
faith of the common people is rejuvenated and strengthened by it; the
commanding words work on the mind uplifting like voices from on high,
refreshing like breaths of divine majesty. But, remarkably enough, such
preaching makes as good as no impression on the man of culture. Taking
his stand exclusively on the standpoint of humanity, he regards the whole
machinery of supernatural driving wheels as an apparatus of motion that
human imagination, without rightly knowing it, has itself constructed.
The effect the machinery has on believing souls is naturally an object of
his observing attention; but it is only an object. It lightens from the
gloom of the clouds, from the dark wings on which the spirit hovers, but
the lightnings do not strike the humane \opage{100}observer; he stands calmly under
the thunder of Sinai, unperturbed, like a correspondent reporting to his
newspaper.

And what, then, comes to stand in the newspaper? ``It is a sign of the
times, a remarkable phenomenon, to see the eloquent N.~N., glowing with
inspiration, gather the masses in crowds about his pulpit. A grand
awakening! His speech is a cloudburst; the words fall like raindrops,
refreshing after a long drought.'' Then follows the description of a
remarkable man with a high forehead, sparkling eyes, sunken cheeks, thin
lips and a long beard. Looks, gestures, turns of phrase and catchwords
are pointed out most exactly, but of the spiritual in the matter itself
not a word is reported.

That God himself should have dictated the ten commandments of the law of
Moses will be found, on behalf of humanity, still more unreasonable than
that he should have dictated how many and how long curtains were to be
hung in the tabernacle, of what colors and with how many loops. The
humanists, who see so clearly that there is a difference between moral
commandments and crimson curtains, also see that the difference consists
precisely in this: that the moral commandments must have what curtains
can hardly have, an unmistakable stamp of the necessity of reason, of
universal validity. It might indeed, on a preliminary consideration, seem
as if it were less needful that God himself should give revelations about
the curtains than that he should reveal the moral laws; but from a
humanist point of view the matter appears the other way round. About
veils, curtains, lamp oil and the like Jehovah necessarily had to give
word-for-word instructions if he wanted the tabernacle arranged exactly
after his own head, for such specialities, which cannot possibly be
derived from pure reason, must for that very reason \opage{101}be reckoned to the
inventory of ``the Kingdom of God.'' As regards the moral laws, on the
other hand, a miraculous revelation is altogether superfluous, for these
laws have their sufficient ground in reason and in human nature, and
belong expressly to the kingdom of humanity.

So sharply can the exclusively religious and the exclusively humane
stand against each other. But that the religious, by one-sidedly holding
fast to the supernatural, just as surely misunderstand religion as the
humane, by one-sidedly keeping to the natural, misunderstand humanity,
can be clearly shown.

Religious awakenings have at all times the great significance that the
truly awakened will grasp the matter from the ground up, put the
original into force and penetrate to the core of the life of the spirit.
But the limit and barrier is that, with all their zeal for the right of
the spirit and the holy laws, they never get as far as any real account
of the natural laws of actual life, of men's laboring, working and
striving in the present. They constantly stop at the explanation that it
is an evil, ungodly, sinful, vain and corruptible world in which we live.
They can count off on their fingers the many transgressions, sins and
dissipations to which the children of the world are given. But what do
they put in their place? A system of holiness that consists in making
life as sour for oneself as possible. It is not enough to keep the ten
commandments; there are many other commandments to keep: God's children
must not, e.g., laugh, smoke tobacco, dance, play cards, go to the
theater, etc. By banishing the enjoyments of life and setting narrow
limits to what is permitted, one aims at securing for the spirit a
lasting dominion over \opage{102}the flesh; and it is certain that, if it were only
a question of a life with fewer and fewer pleasures, the brothers and
sisters of the strict discipline would stand on a superior level. But the
false presupposition that spirit and nature stand hostile to each other
shines through. So long as it has not been seen on the religious side
that the spiritual needs the natural, that it is the goal of the spirit
to refine the natural, to develop, form and further human life in all
directions, so long will the proclamation of the Kingdom of God be
powerless in the face of humanity, without essential influence on the
existence of the present age as a whole.

At the stage of humanity the trouble is of the opposite kind. Here the
start is made without further ado with the natural, the present, the
manifold tasks of actual life. Then the point is to show the laws. But
what sort of laws are they? Sheer natural laws---valid for the animal as
well as for man. In human life these laws are now, so far as the activity
of the brain permits, to be transformed by degrees into general laws of
reason and thereby into laws of the spirit. But the transformation goes
only slowly. Pure reason with its categorical imperative is an
abstraction, and nature rebels against all abstractions. Actual reason
lies hidden in the actual drives. The actual drives are not only selfish
but also sympathetic. In the mutual helpfulness of the baboons the moral
essence of humanity is prefigured. In man---especially through the
hereditary improvement of the brain---empirical ideas become little by
little a priori. Under the continued interaction of the manifold forces of
society, the moral drives pass imperceptibly over into moral \opage{103}principles,
into a system of moral laws. The moral laws form the foundation of
humanity and the humane spirit. Trade, industry, economy, capital,
politics, science, poetry and art labor in the service of the spirit; the
result of all these labors is the growing power of the human spirit over
lower nature. This spirit of humanity has an infinitely long past and an
infinitely long future, but, alas, no actual present; for it has never
had and will never get any actual unity of will. From the interaction
between the single wills of society the unity of will, and with it actual
spirit, can never result. The single wills are, in their inmost ground,
wills of drive; but when the will of drive is to be the original, then we
come---to the baboons.

The exclusively religious go wrong in the absolute; they stop at a
supernatural will of God whose unconditional commandments they are
unable, in all their obedience, to conceive as proclamations of a will of
the spirit that permeates the whole of human life. They cannot satisfy
religion because they cannot satisfy humanity. The exclusively humane go
wrong in the relative. With their natural wills of drive they are not
able to raise the human spirit to actual dominion over the conflicting
elements of natural life, and, as a matter of course, not able to satisfy
the demands of humanity. Religion demands the humane, for it demands a
divine liberation of the powers of human life. Humanity requires the
religious, for the exertions of relative wills are in vain if they do not
have the holy, absolute will as their bearer and prime mover.

\opage{104}\emph{Justice}. ``Tremble, you sinners, for God is just!'' When this cry
rings out over the world in the name of revelation, it seldom fails of its
effect. Just when men are most at odds, they are most of all agreed in
complaining of the abomination of wrong, the presumption, violence and
arbitrariness of growing injustice. On the religious feeling of the
common people that admonishing cry will then work like a stroke of
lightning and give consciences a mighty jolt. The believers beat their
breasts; everyone confesses his sin in order to flee from the wrath to
come. Behind this momentary penitence there lurks not infrequently a
selfish notion that by confessing one's sins one can at once be rid of
them, that it is of course not so much oneself as the ungodly world, full
of injustice, that must dread the punishment. The judge of the world, the
Just One, takes up the cause of the oppressed, the wronged, and everyone,
after all, feels himself wronged. ``When the Son of Man comes in the
clouds with power and great glory, he shall separate the just from the
unjust as a shepherd separates the sheep from the goats, and then all the
kindreds of the earth shall wail; but the elect, who are atoned for by the
blood of the Lamb, who have put off their own righteousness and put on
the Atoner's, shall be received into grace.''

On the man of culture of the present age such a solemn announcement of
God's justice will make no impression. Criticism must find that the
sharp separation between the just and the unjust is inhumane: there is
right and wrong on both sides, on that of the ``sheep'' as well as that
of the ``goats.'' \opage{105}Nor is it correct to say: God is just, for justice is a
concept that finds application only to human relations. Even if one will
say that justice is divine, it is nonetheless only an abstraction. Pure
justice is worth nothing in itself; it is a commemorative coin that is
not in circulation. Actual justice exists only in the actual disposition
that shows itself in the strife about mine and thine, and so in the way in
which one conceives and treats positive questions of right. But the
positive questions of right are human questions, and the right of pardon
a human right with a smack of arbitrariness, which characterizes finitude
and belongs in each and every respect to relativity.

From this standpoint, too, the opposition between the religious and the
humane shows itself as an irreconcilable discord. But what is the reason?
In the immediately religious notion it looks as if human justice were
without validity before God. It can then be said of all appearances of
this justice what a Father of the Church has said of the virtues of the
heathen: that they are splendid vices. If this is understood to mean that
human justice, when, held fast for itself, it will have merit before God,
must become an expression of human effrontery, nothing can be objected to
it. But when justice is conceived in such a way that this essence must
necessarily exclude human nature and put the divine in its place, then
the religious is indeed absolutely incompatible with the humane, but
then the fundamental religious thought is also obscured. Believers are to
appropriate the Atoner's righteousness; but the Atoner makes it manifest
precisely that \opage{106}the divine itself is human. If the Atoner is conceived
merely as a historical individuality who once suffered and died because
of men's injustice, then it is altogether impossible to appropriate his
righteousness. But the essence of the Atoner is, in the divine-human
sense, the inmost true essence of mankind. The Atoner came once for all
externally to men as a historical figure precisely in order that he may
now and henceforth come inwardly in the spirit with his righteousness,
not as a foreign person but as the essence and ground in all true human
personality. To appropriate the Atoner's righteousness is to deify human
justice and to humanize divine justice; it is, in the deepest sense, to
satisfy humanity.

If we take our stand on the humane standpoint and begin with the strife
about mine and thine, then it is granted, after all, that disputes about
rights, even the most zealous, cannot possibly lead to just relations
unless they are grounded on the thought of justice and conducted with a
just disposition and a just will. If one does not respect the
unconditional validity of justice, but makes justice in and for itself an
empty concept, then the finite strife about rights will end with the
right of the stronger, and relative justice will be transformed into a
positive expression of relative injustice. If relative justice is to
triumph approximately over injustice, it must happen by virtue of the
idea of justice. But the idea of justice is an idea of will; absolute
justice presupposes an absolutely just will, a will that in the last
instance gets the better of all injustice. Faith in triumphant justice is
inseparable from faith in the God of justice. \opage{107}For the sake of human
justice God must be absolutely just: the humane requires the religious.

What is here indicated can be developed infinitely; but the further the
development goes, and the more exactly the determinations of the relation
are stated, the more clearly it appears that religion must be humane, and
humanity religious.

\emph{Spiritual formation}. If we set the existence of the present age, on
its historical foundation, with its labors and ends, freely over against
the apostolic existence with its ideal breakthroughs, while yet
recognizing both existences as essentially belonging together, then we
get an impression of a dualism in the grand style. The servants of the
Kingdom of God, with their heavenward direction of life and all-sacrificing
renunciation of the world, ``by honor and dishonor, by evil report and
good report, as deceivers and yet true, as unknown and yet well known, as
dying and yet living, as chastened and yet not killed, as sorrowful and
yet rejoicing, as poor yet making many rich, as having nothing and yet
possessing all things''---and then the servants of humanity, immersed in
the earthly, striving for the temporal, laboring for the present, for the
demands of the moment in the bondage of finitude, rich and yet poor,
surfeited and yet hungry, cheerful and yet in secret anxiety, glad and yet
troubled, always troubled about worldly goods, honor, power and civic
influence---what an enormous opposition!

And yet both existences belong to one and the same realm, the kingdom of
the spirit; the fundamental law is the same for both, the fundamental law
of freedom, justice and love. The unconditional demand is that the spirit
shall rule everywhere in its kingdom, in the lower as \opage{108}in the higher
regions. Now if the apostolic existence did not stand firm on the holy
ground of authority, the testimony to the right and the goal of the
spirit would vanish. What the spirit can demand of man when sensuous drive
and selfish willfulness rebel against its dominion, what sacrifices it
requires, what powers it has, what means it offers when in the strife it
comes to the uttermost---all this would be clean forgotten; the holy, the
absolutely exalted, which holds the standard, would sink down into sheer
finitude, sheer relativity, and false weight and measure would creep in
everywhere and confuse humane relations. As little as right and justice
can be done in a civil society with false coin, false measure and weight,
so little can order be kept in the kingdom of the spirit when the ideal
standards are broken, when the higher can become the lower and the lower
the higher, as the day by turns raises up and casts down its idols.

The ideal existence holds not only for God's sake but first and last for
man's sake, and the present age is so far right when it fixes its gaze on
the world of realities and holds to the humane. It is undeniably an
illusion, indeed an ingrown prejudice, that the holy God of the spirit
should have prescribed to men a single law, a single condition, that was
not grounded in the disposition of human nature, not necessary for its
free, personal self-development. But that the personal can be grounded
on something impersonal is an illusion too.

It is the life of the congregation that continues to form the foundation
of the life of the spirit. The servants of the word in the congregation
have the great task of upholding and enjoining the conditions of the
Kingdom of God, not with an assumed apostolic \opage{109}authority but, for the sake
of appropriation, with the whole superiority of spiritual formation. The
teachers of youth have the great task, by the means of spiritual
formation, of giving the revealed truths a practically intelligible
expression in the language of humanity, and of making them present, by
the authority of the ideals, in the earnestness of upbringing.---Then
there is indeed a difference between the church-school and the
culture-school, but the discord is removed, a mutual understanding has
been reached, so that both schools, if by different means and along
different ways, can work in the same spirit toward the attainment of the
same goal.

There are, to be sure, those who think that the time of religion is past,
that humanity alone has the future before it. On their instructions we
are not to remain standing on the ground of the Church and stare toward
the East with Luther; we are, with a new Columbus, to hoist the swaying
flag of ``free thought'' and boldly steer our course toward the West. ``Far
in the West, in the Atlantic of the discoveries of nature, lies the
promised land of the spirit, hidden like a pearl in the deep.''---That is
false guidance. He who blindly plunges into the gulf of naturalism does
not come up ``with a pearl in his hand,'' but---with a bubble in his
mouth.

The same spirit of life and truth that once, by extraordinary
self-communication, brought forth the apostolic existence still
communicates itself, on the humane conditions of appropriation, with the
power of rejuvenation and transformation to the existence of the present
age. The scheme is different, but the content is the same. The means we
have; the conditions are given us; it depends only on our using them
rightly. To use them rightly, upbringing and spiritual formation are
required.

\end{document}
